% ---------------------------------------------------------------------------
% Supplementary Material  (anonymized for double-blind submission: the system
% name and its source package are redacted; restore both in the camera-ready)
%
% STANDALONE document.  It deliberately does NOT use aaai2027.sty (which is
% submission-locked) and does not \input any file from sections/.  Build with:
%     latexmk -pdf -interaction=nonstopmode supplement.tex
% It resolves citations against ../references.bib.
% ---------------------------------------------------------------------------
\documentclass[11pt]{article}

\usepackage[letterpaper,margin=1in]{geometry}
\usepackage{booktabs}
\usepackage{amsmath}
\usepackage{amssymb}
\usepackage{amsthm}
\usepackage[numbers]{natbib}
\usepackage{array}
\usepackage{longtable}
\usepackage{graphicx}
\usepackage{xcolor}
\usepackage[T1]{fontenc}
\usepackage{lmodern}
\AtBeginDocument{\DeclareFontShape{T1}{lmr}{bx}{sc}{<-> ec-lmcsc10}{}}
\usepackage[utf8]{inputenc}
\usepackage[hidelinks]{hyperref}

\newtheorem{theorem}{Theorem}
\newtheorem{lemma}{Lemma}
\newtheorem{definition}{Definition}
\newtheorem{assumption}{Assumption}
\newtheorem{proposition}{Proposition}
\newtheorem{remark}{Remark}

\newcommand{\Traces}{\mathrm{Traces}}
\newcommand{\Tev}{\Traces_{\mathrm{ev}}}
\newcommand{\Lang}{\mathcal{L}}
% Literal artifact paths and corpus slugs are long and contain no natural break
% points, so we make "_" (and "/" via \sfrac-free \allowbreak below) breakable
% inside the monospace runs.  Purely typographic; the text is unchanged.
\DeclareRobustCommand{\bslash}{\char`\_\allowbreak}
\DeclareRobustCommand{\artifact}[1]{{\small\ttfamily\let\_\bslash #1}}
\DeclareRobustCommand{\slug}[1]{{\ttfamily\let\_\bslash #1}}
\sloppy

\IfFileExists{../generated/human_agreement.tex}{%
  % Generated by scripts/human_agreement.py -- do not edit by hand.
\newcommand{\humanNflags}{0}
\newcommand{\humanKappa}{--}
\newcommand{\humanAlpha}{--}
\newcommand{\humanRaw}{--}
\newcommand{\humanLLMkappaA}{--}
\newcommand{\humanLLMkappaB}{--}
\newcommand{\humanGraphN}{0}
\newcommand{\humanGraphEdgeF}{--}
\newcommand{\humanGraphKind}{--}
\newcommand{\humanRefEdgeF}{--}
%
}{%
  \newcommand{\humanNflags}{0}
  \newcommand{\humanKappa}{--}
  \newcommand{\humanAlpha}{--}
  \newcommand{\humanRaw}{--}
  \newcommand{\humanLLMkappaA}{--}
  \newcommand{\humanLLMkappaB}{--}
  \newcommand{\humanGraphN}{0}
  \newcommand{\humanGraphEdgeF}{--}
  \newcommand{\humanGraphKind}{--}
  \newcommand{\humanRefEdgeF}{--}
}
\ifnum\humanNflags>0
  \newcommand{\humanValidationStatus}{human-validated}
\else
  \newcommand{\humanValidationStatus}{human validation pending}
\fi

\title{Supplementary Material for\\
\emph{Before Verifying Agent Workflows: Auditing Extraction Fidelity and Policy
Applicability}}
\author{}
\date{}

\begin{document}
\maketitle

\begin{abstract}
\noindent
This supplement discharges the forward references made by the main paper, carries
the positioning table displaced from it by the page limit, and adds three
reviewer-facing methods documents that the main text has no room for. Every
numeric claim below is annotated with the committed artifact path it
comes from, or with the script that was executed to produce it. Where a promised
item cannot be produced honestly from the artifacts, we say so explicitly and
explain why rather than substituting a number of unclear provenance;
Sections~\ref{sec:ontology} and~\ref{sec:stubs} each contain one such
disclosure, and Section~\ref{sec:proof} states its central result as an
explicitly \emph{conditional} theorem whose hypothesis the implementation does
not verify.
\end{abstract}

\tableofcontents

\bigskip
\hrule
\bigskip

\noindent\textbf{Artifact conventions.} Paths are relative to the repository
root. \artifact{corpus/real\_world/} holds the mining study's committed JSON;
\artifact{scripts/} holds the executable analyses. The submitted
\artifact{reproduce\_all.sh} regenerates HGP-1 before every downstream summary,
rebuilds matched-set and confidence outputs, and finishes with
\artifact{validate\_submission\_claims.py}, which fails if any primary artifact
reverts to the withdrawn $3/119$ policy union, the pre-exclusion 187-flag set,
or the collision-contaminated fidelity set.

% ===========================================================================
\section{The Orthogonal Attribute Ontology and its Per-Framework Mapping}
\label{sec:ontology}

\emph{Promised at} \artifact{sections/02\_system.tex}, paragraph
\emph{An orthogonal attribute ontology}: ``The full vocabulary and its
per-framework mapping appear in the supplement.''

\subsection{The vocabulary as implemented}

The vocabulary below is transcribed from
\artifact{src/[pkg]/graph/model.py} (the frozensets \slug{VALID\_EFFECTS},
\slug{VALID\_CAPABILITIES}, and the \slug{back\_edge} field of
\slug{GraphEdge}), not from the paper's prose, so the two are guaranteed to
agree.

\begin{table}[htbp]
\centering\small
\caption{The \emph{effect} set: observable side effects a node's execution may
have on the world or on shared state. Non-exclusive---a node may carry any
subset, including the empty set. Source: \artifact{src/[pkg]/graph/model.py},
\slug{VALID\_EFFECTS}.}
\label{tab:effects}
\begin{tabular}{@{}lp{0.62\linewidth}@{}}
\toprule
\textbf{Effect} & \textbf{Definition (verbatim from the implementation)} \\
\midrule
\slug{read}        & reads external data / state (queries, retrieval, fetch) \\
\slug{write}       & writes/persists data (db insert, file write, state store) \\
\slug{execute}     & runs code / commands (code exec, shell, kubectl, restart) \\
\slug{communicate} & sends messages outward (email, chat, API notify, publish) \\
\slug{financial}   & moves money or places trades (trading/banking/payment) \\
\slug{delete}      & destroys/removes data or resources \\
\bottomrule
\end{tabular}
\end{table}

\begin{table}[htbp]
\centering\small
\caption{The \emph{capability} set: structural abilities a node possesses,
orthogonal to \slug{NodeKind}. A single node may hold several; the
implementation's own worked example is an LLM node that also invokes tools and
routes, i.e.\ \slug{("llm","invokes\_tool","routes")}. Source:
\artifact{src/[pkg]/graph/model.py}, \slug{VALID\_CAPABILITIES}.}
\label{tab:capabilities}
\begin{tabular}{@{}lp{0.62\linewidth}@{}}
\toprule
\textbf{Capability} & \textbf{Definition (verbatim from the implementation)} \\
\midrule
\slug{llm}          & performs an LLM/model inference call \\
\slug{invokes\_tool} & invokes one or more bound tools \\
\slug{human\_pause}  & pauses for human input / approval \\
\slug{routes}       & makes a branching/routing decision over successors \\
\slug{state\_update} & mutates shared graph state \\
\slug{subgraph}     & delegates to / expands a nested subgraph \\
\bottomrule
\end{tabular}
\end{table}

\paragraph{The edge back-edge flag.}
\slug{GraphEdge} carries \slug{back\_edge: bool = False} alongside
\slug{kind}. \slug{EdgeKind.LOOP} is retained for backward compatibility,
but the orthogonal flag is the credible encoding: a conditional loop is
(\slug, \slug{back\_edge=True}), which the single
\slug{EdgeKind} label cannot express. Back edges are computed in
\artifact{scripts/build\_corpus.py} (\slug{\_back\_edges}) by a classical
DFS---an edge whose target is on the recursion stack---so the flag is exactly
the set of loop-closing edges, not a name heuristic.

\paragraph{Relationship to the exclusive kinds.}
The single-label vocabularies remain in the model as the compatibility view and
are what every check still keys on:
\slug{NodeKind} $\in$ \{\textsc{entry}, \textsc{exit}, \textsc{tool},
\textsc{llm}, \textsc{router}, \textsc{human}, \textsc{subgraph},
\textsc{passthrough}\} and \slug{EdgeKind} $\in$ \{\textsc{direct},
\textsc{conditional}, \textsc{parallel}, \textsc{loop}\}. Note that
\textsc{passthrough} is a first-class \slug{NodeKind} in the implementation
even though earlier drafts of the paper's formal vocabulary omitted it.

\subsection{Per-framework mapping}

\begin{table}[htbp]
\centering\small
\caption{Per-framework mapping into the exclusive \slug{NodeKind}/%
\slug{EdgeKind} view, transcribed from the four runtime extractors under
\artifact{src/[pkg]/graph/extract/}. ``sy'' marks synthesized sentinels
(elements fabricated by the extractor, provenance
\slug{origin=synthesized}).}
\label{tab:fwmap}
\begin{tabular}{@{}lp{0.70\linewidth}@{}}
\toprule
\textbf{Framework} & \textbf{Mapping} \\
\midrule
LangGraph
& \slug{START}$\to$\textsc{entry}; \slug{END}$\to$\textsc{exit};
  \slug{ToolNode} / declared tool binding$\to$\textsc{tool};
  interrupt/human-named node$\to$\textsc{human} (name heuristic, so
  \slug{confidence=may}); conditional-branch source$\to$\textsc{router}
  (heuristic); otherwise \textsc{llm}. Edges from
  \slug{add\_conditional\_edges}$\to$\textsc{conditional}; from
  \slug{add\_edge}$\to$\textsc{direct}. Nodes and \slug{add\_edge} edges are
  \slug{origin=runtime}, \slug{confidence=exact} unless the kind was guessed.
  \\[3pt]
CrewAI
& \slug{Task} with bound tools$\to$\textsc{tool}, else \textsc{llm};
  hierarchical-process manager$\to$\textsc{router} (sy);
  \slug{\_\_entry\_\_}/\slug{\_\_exit\_\_} sentinels (sy). Task \emph{ordering}
  is inferred from the crew's task list, so the chaining edges are
  \slug{origin=runtime}, \slug{confidence=may}; manager fan-out edges are
  \textsc{conditional} and synthesized. \\[3pt]
AutoGen
& \slug{UserProxyAgent}-like participant$\to$\textsc{human};
  \slug{SelectorGroupChat} selector$\to$\textsc{router}; other participants
  $\to$\textsc{llm}; \slug{\_\_entry\_\_}/\slug{\_\_exit\_\_} (sy).
  Participants are \slug{confidence=exact}; the transition relation
  reconstructed from team semantics (round-robin ring, selector fan-out) keeps
  \slug{origin=runtime} with a weaker confidence, and hook-up edges to the
  sentinels are \slug{confidence=heuristic}. \\[3pt]
Google ADK
& \slug{SequentialAgent}/\slug{ParallelAgent}/\slug{LoopAgent}
  $\to$\textsc{subgraph} (\slug{confidence=exact}) with children expanded;
  \slug{LongRunningFunctionTool}/confirmation callback$\to$\textsc{human};
  agent with declared tools$\to$\textsc{tool}; else \textsc{llm}.
  \slug{SequentialAgent} children are chained by \textsc{direct} edges,
  \slug{ParallelAgent} children by \textsc{parallel} edges, and
  \slug{LoopAgent} closes an \textsc{edgekind.loop} edge from the last child
  back to the first (or a self-loop for a single child). All composition edges
  are \slug{origin=runtime}, \slug{confidence=may}, because the ordering is
  inferred from the composite's semantics rather than declared edge-by-edge. \\
\bottomrule
\end{tabular}
\end{table}

\subsection{Disclosure: the effect/capability layer is not populated by the extractors}
\label{sec:ontology-gap}

We must be precise about what is implemented, because the main text's phrase
``its per-framework mapping'' could be read as a claim we cannot support.

A repository-wide search for assignments to the \slug{effects} and
\slug{capabilities} fields outside the model definition returns exactly two
non-test sites, both in \artifact{scripts/build\_corpus.py}. \textbf{None of the
four per-framework extractors under \artifact{src/[pkg]/graph/extract/}
populates \slug{effects} or \slug{capabilities}, and none sets
\slug{back\_edge}.} Every graph produced by extraction---runtime or AST,
curated or mined---therefore carries the dataclass defaults \slug{effects=()},
\slug{capabilities=()}, \slug{back\_edge=False}. The layer is, as the model's
own comment puts it, ``a purely additive descriptive layer'' that ``does not
disturb \slug{NodeKind}/\slug{EdgeKind} or any check logic''.

The only mapping into the ontology that actually runs is the one applied when
the \emph{curated} corpus is authored (\artifact{scripts/build\_corpus.py},
functions \slug{\_tag\_node} and \slug{\_effects\_for}), and it is
framework-independent:

\begin{itemize}
\item \textbf{Capabilities} are a deterministic function of \slug{NodeKind}
alone (\slug{\_CAPABILITY\_BY\_KIND}):
\textsc{llm}$\to$\slug{("llm",)},
\textsc{tool}$\to$\slug{("invokes\_tool",)},
\textsc{router}$\to$\slug{("routes",)},
\textsc{human}$\to$\slug{("human\_pause",)},
\textsc{subgraph}$\to$\slug{("subgraph",)}, and
\textsc{entry}/\textsc{exit}/\textsc{passthrough}$\to$\slug{()}. Consequently
\slug{state\_update} is in the vocabulary but is \emph{never assigned} by any
code path in the repository; it is currently vocabulary-only.
\item \textbf{Effects} are a keyword match over a node's \emph{declared tool
names} (\slug{\_EFFECT\_KEYWORDS}, 30 keyword rules), e.g.\
\slug{delete}/\slug{remove}/\slug{purge}/\slug{drop}$\to$\slug{delete};
\slug{payment}/\slug{trading}$\to$\slug{financial} (with
\slug{banking}$\to$\slug{financial,write});
\slug{smtp}/\slug{email}/\slug{notify}/\slug{social\_media}$\to$%
\slug{communicate};
\slug{execute}/\slug{pytest}/\slug{kubectl}/\slug{restart}/%
\slug{rollback}$\to$\slug{execute};
\slug{insert}/\slug{dall}/\slug{canva}$\to$\slug{write} (with
\slug{cms}$\to$\slug{write,communicate}); and a read block
(\slug{search}, \slug{query}, \slug{fetch}, \slug{retriev},
\slug{scrap}, \slug{parse}, \slug{vision}, \slug{finance},
\slug{bloomberg}, \slug{brokerage}, \slug{market\_data}, \slug{bureau},
\slug{verify}, \slug{analytics}, \slug{survey}, \slug{competitor},
\slug{scan}, \slug{validator})$\to$\slug{read}. Genuinely ambiguous tool
names are deliberately left with no effect.
\end{itemize}

\begin{remark}[Consequence for the paper's claims]
Because the ontology is populated only on authored graphs, and only from
declared tool names, it inherits exactly the blindness the main text diagnoses:
a node that invokes a side-effecting tool inside its body with no declared
binding receives the empty effect set. This is the same
\slug{tool}$\rightarrow$\slug{llm} kind error that dominates LangGraph's
residual extraction error, viewed through a different field. \textbf{A genuine
per-framework mapping of framework constructs onto effects and capabilities
during extraction is not implemented and is future work.} Table~\ref{tab:fwmap}
documents the per-framework mapping that \emph{is} implemented---the one into
the exclusive kinds---and should be read as such.
\end{remark}

% ===========================================================================
\section{Event-Trace Soundness: A Conditional Theorem}
\label{sec:proof}

\emph{Promised at} \artifact{sections/02\_system.tex}, paragraph
\emph{Soundness requires an event-complete abstraction}: ``the full proof is in the
supplement.''

\subsection{Why the existing graph-path lemmas do not suffice}

An earlier (arXiv) appendix proves eight lemmas, of which the relevant one
(\emph{Static Temporal Soundness}) establishes: if the graph\,$\times$\,DFA
product reports no violation, then no \emph{path through $G$} violates the
property. The main text now says plainly that this is not the guarantee we need.
The gap is not a technicality. A graph path is a sequence of \emph{nodes}; a
policy is evaluated over a sequence of \emph{events}. A node whose body issues a
tool call that appears nowhere in its declared bindings contributes an event that
the graph language does not contain, so a graph that over-approximates every
path can still under-approximate the event language. Path containment and event
containment are logically independent, and the machinery only ever establishes
the former. This section therefore states and proves the event-level result
directly, and is explicit that its principal hypothesis is \emph{assumed},
never checked.

\subsection{Setting}

\begin{definition}[Event alphabet and concrete traces]
Fix a finite event alphabet $\Sigma$. An \emph{event} is a tool invocation with
its identity (and, where the policy inspects them, its arguments), a
node-emitted event, or the distinguished termination marker. For a workflow
program $P$, let $\Tev(P)\subseteq\Sigma^{*}$ be the set of finite event traces
emitted by \emph{terminating} executions of $P$, under whatever scheduler,
model, and environment $P$ is deployed with. Non-terminating executions
contribute no member of $\Tev(P)$.
\end{definition}

\begin{definition}[Graph event language]
Let $G=(V,E,v_0,V_f)$ be the extracted graph and let
$\lambda: V \to 2^{\Sigma^{*}}$ assign to each node the set of finite event
sequences the abstraction permits that node to emit. For a tool node $v$ with
declared binding set $T(v)$, the checker takes
$\lambda(v) = T(v)^{*}$---every finite invocation sequence over the declared
tools, \emph{including the empty sequence}---which is the conservative closure
described in the main text. For a non-emitting node, $\lambda(v)=\{\varepsilon\}$.
Let $\Pi(G)$ be the set of \emph{maximal} finite paths of $G$ from $v_0$: paths
ending at a node in $V_f$ or at a node with no successor. Then
\[
  \Lang(G,\lambda) \;=\;
  \bigl\{\, w_0 w_1 \cdots w_k \;\bigm|\;
    v_0 v_1 \cdots v_k \in \Pi(G),\
    w_i \in \lambda(v_i)\text{ for }0\leq i\leq k \,\bigr\}
  \;\subseteq\; \Sigma^{*}.
\]
\end{definition}

\begin{definition}[Policy and its DFA]
A policy $\varphi$ is an LTL$_f$ formula in the DSL fragment of the main text.
Its compiled deterministic automaton is
$\mathcal{A}_\varphi=(Q,\Sigma,q_0,\delta,F)$ with accepting set $F\subseteq Q$
and violation set $Q_{\mathrm{viol}}\subseteq Q\setminus F$ defined as the
\emph{doomed} states, i.e.\ those from which no state in $F$ is reachable.
A finite trace $\sigma\in\Sigma^{*}$ is \emph{DFA-accepted} iff
$\hat\delta(q_0,\sigma)\in F$; by construction of the LTL$_f$-to-DFA
compilation, $\sigma$ is DFA-accepted iff $\sigma\models\varphi$ in the finite-trace
semantics of \citet{degiacomo2013ltlf}.
\end{definition}

\begin{assumption}[Event-trace conservatism]
\label{as:etc}
$\Tev(P)\subseteq\Lang(G,\lambda)$.
\end{assumption}

\begin{assumption}[Compiler correctness]
\label{as:compiler}
For every $\sigma\in\Sigma^{*}$: $\hat\delta(q_0,\sigma)\in F$ iff
$\sigma\models\varphi$.
\end{assumption}

\subsection{The theorem}

\begin{theorem}[Conditional event-trace soundness]
\label{thm:etc}
Let $P$ be a workflow, $G$ its extracted graph with labelling $\lambda$, and
$\varphi$ a policy with DFA $\mathcal{A}_\varphi$. Suppose
\begin{enumerate}
\item[(H1)] Assumption~\ref{as:etc} holds: $\Tev(P)\subseteq\Lang(G,\lambda)$; and
\item[(H2)] every $\sigma \in \Lang(G,\lambda)$ is DFA-accepted, i.e.\
$\hat\delta(q_0,\sigma)\in F$.
\end{enumerate}
Then every concrete finite terminating execution of $P$ satisfies $\varphi$.
\end{theorem}

\begin{proof}
Let $e$ be any concrete finite terminating execution of $P$ and let
$\sigma_e\in\Sigma^{*}$ be its event trace. By definition of $\Tev(P)$ we have
$\sigma_e\in\Tev(P)$. By (H1), $\sigma_e\in\Lang(G,\lambda)$. By (H2),
$\hat\delta(q_0,\sigma_e)\in F$, so $\sigma_e$ is DFA-accepted. By
Assumption~\ref{as:compiler}, $\sigma_e\models\varphi$. As $e$ was arbitrary,
every concrete finite terminating execution of $P$ satisfies $\varphi$.
\end{proof}

The proof is deliberately trivial. That is the point: once event-trace
conservatism is granted, soundness is immediate, and \emph{all} of the technical
content of the guarantee sits in (H1). The remainder of this section is about
what the implementation does and does not do towards (H1) and (H2).

\subsection{What the implementation discharges: (H2)}

(H2) is the hypothesis the product construction actually establishes, and it
does so by a standard argument, which we give for completeness.

\begin{lemma}[The product decides (H2)]
\label{lem:product}
Define the product state space $V\times Q$ with initial state $(v_0,q_0)$ and
transition relation
\[
  (v,q)\;\longrightarrow\;(v',q')
  \quad\text{whenever}\quad
  (v,v')\in E \ \text{ and }\ q'=\hat\delta(q,w)\ \text{ for some } w\in\lambda(v).
\]
Let $R$ be the set of product states reachable from $(v_0,q_0)$. If
(i) no $(v,q)\in R$ has $q\in Q_{\mathrm{viol}}$, and
(ii) for every $(v,q)\in R$ with $v$ a termination point (i.e.\ $v\in V_f$ or $v$
has no successor) and every $w\in\lambda(v)$ we have $\hat\delta(q,w)\in F$,
then (H2) holds.
\end{lemma}

\begin{proof}
Take $\sigma\in\Lang(G,\lambda)$. By definition
$\sigma=w_0\cdots w_k$ for some maximal path
$v_0v_1\cdots v_k\in\Pi(G)$ and choices $w_i\in\lambda(v_i)$. Define
$q_i=\hat\delta(q_0,w_0\cdots w_{i-1})$ for $0\leq i\leq k$, where the empty
prefix gives $q_0$. An induction on $i$ shows $(v_i,q_i)\in R$: the base case is
$(v_0,q_0)\in R$, and the step from $i$ to $i+1$ follows
$(v_i,v_{i+1})\in E$ while consuming $w_i\in\lambda(v_i)$, hence reaches
$(v_{i+1},q_{i+1})$. Now $v_k$ is a termination point because the path is
maximal, so condition (ii), applied to $w_k\in\lambda(v_k)$, gives
$\hat\delta(q_k,w_k)=\hat\delta(q_0,\sigma)\in F$. Since $\sigma$ was
arbitrary, (H2) holds. Condition (i) is not needed for this
implication; it is the earlier-warning mechanism---a reachable doomed state
witnesses a bad prefix, and the checker reports \slug{may\_violate} at the
first event that makes every extension violating rather than waiting for
termination. The two mechanisms are disjoint by construction, since
$Q_{\mathrm{viol}}\cap F=\emptyset$.
\end{proof}

\begin{remark}[Why $\lambda(v)=T(v)^{*}$ and not $T(v)$]
The closure over \emph{all} finite invocation sequences, including
$\varepsilon$, is what makes (H2) robust to the two things static extraction
cannot decide about a tool node: \emph{which} of the declared tools fires, and
\emph{whether any} fires. Restricting $\lambda(v)$ to single invocations would
make $\Lang(G,\lambda)$ smaller and (H2) easier to satisfy---and unsound, since
a real execution invoking a declared tool twice would fall outside it. The
closure is the reason a false alarm is possible and a false proof is not,
\emph{given} (H1).
\end{remark}

\begin{remark}[Divergent cycles]
Lemma~\ref{lem:product} quantifies over \emph{maximal finite} paths. If the
product has a reachable cycle on which an $\mathbf{F}$-obligation stays pending,
the checker returns \slug{inconclusive} rather than enumerating the
terminating traces that thread the cycle. Universal satisfaction over the finite
terminating traces of a finite graph is decidable, so this is an
engineering-scope choice and not a semantic limitation; it makes the checker
weaker, never unsound.
\end{remark}

\subsection{What the implementation does \emph{not} discharge: (H1)}
\label{sec:h1}

\begin{proposition}[Exact provenance does not entail (H1)]
\label{prop:gap}
Let \slug{region\_is\_provenance\_exact} mean that every traversed edge and
every visited node carries \slug{exact} provenance. Then
\slug{region\_is\_provenance\_exact}$(G)$ does not imply
$\Tev(P)\subseteq\Lang(G,\lambda)$.
\end{proposition}

\begin{proof}[Justification]
Exact provenance is a statement about \emph{how each recorded element was
obtained}---this edge was read from a framework declaration rather than
guessed---and quantifies only over elements that were recorded. It says nothing
about elements that were never recorded. Concretely: let $v$ be a node whose
body calls a tool $t$ directly, with no declared binding. Every edge incident to
$v$ may be read exactly from \slug{add\_edge}, and $v$ may have $T(v)=\emptyset$
recorded exactly, so \slug{region\_is\_provenance\_exact} holds; yet
$\lambda(v)=\emptyset^{*}=\{\varepsilon\}$ while the concrete execution emits
$t$, so $\sigma_e\notin\Lang(G,\lambda)$ and (H1) fails. Nothing in
\slug{region\_is\_provenance\_exact} inspects node bodies for unrecorded events.
\end{proof}

This is not a hypothetical counterexample. It is the
\slug{tool}$\rightarrow$\slug{llm} kind error that dominates LangGraph's
residual extraction error in the main paper's fidelity study
(LangGraph node-kind accuracy $0.670$; Table~\ref{tab:fidelity_v2}), and the
same error hides side-effecting tools from the risk-aware gate in 32 of 119
sampled workflows (source: \artifact{corpus/real\_world/gt\_triage\_results.json},
via \artifact{corpus/real\_world/error\_decomposition.json}).

\paragraph{How (H1) is actually obtained.}
In the implementation, (H1) is discharged in exactly one way: \textbf{by
assertion}. The caller passes \slug{assume\_trace\_conservative}, asserting
that the graph is an authored abstraction whose nodes emit no events outside
their declared bindings. Provenance qualifies witnesses but is not a proof of
(H1) and never authorizes \slug{safe} by itself. Accordingly:

\begin{itemize}
\item On the \textbf{curated corpus}, authorship supplies the assertion, and every
pruning number in the main text is sound \emph{relative to that assertion},
which the paper labels as such.
\item On \textbf{mined graphs}, the caller supplies no event-coverage assertion,
so the gate withholds \slug{safe} outright. No false proof is emitted, because
no proof is emitted.
\item In both cases a \slug{may\_violate} finding is never suppressed by the
gate, so the gate can only lose precision, never soundness.
\end{itemize}

\begin{remark}[The missing piece, stated as such]
Establishing (H1) independently requires a body-level effect analysis certifying
that a node emits no events outside its declared bindings---for example, a
sound over-approximation of the callees reachable from each node's body, closed
under the framework's dispatch mechanisms. \textbf{This is not implemented.}
Theorem~\ref{thm:etc} is therefore a conditional result whose antecedent the
system does not verify, and the honest summary of the guarantee is: \emph{the
product is sound relative to an authored abstraction whose event coverage is
asserted by the caller}. We state it this way rather than in the unconditional
form because the unconditional form would be false of the implementation.
\end{remark}

\begin{remark}[Scope of Assumption~\ref{as:compiler}]
Assumption~\ref{as:compiler} is not verified either, but it is
\emph{tested}: the compiler's verdicts are differentially checked against an
independent, directly recursive reference evaluator---exhaustively over all
traces up to length six for two-atom formulas, and property-based beyond. That
is evidence, not proof, and we label it accordingly. Unlike (H1), however, it
concerns a component we wrote and can in principle verify; (H1) concerns
arbitrary third-party Python.
\end{remark}

% ===========================================================================
\section{Per-Framework Fidelity: Clustered CIs and the Unmatched Breakdown}
\label{sec:fidelity}

\emph{Promised at} \artifact{sections/03\_realworld.tex}, caption of
Table~\emph{tab:realworld\_fidelity}: ``Repository-clustered bootstrap CIs and
the unmatched breakdown are in the supplement.''

\subsection{Provenance of these numbers}

Point estimates come from
\artifact{corpus/real\_world/matched\_fidelity.json}, key
\slug{provenance\_collisions.collision\_free\_fidelity} ($n=106$). Intervals
were computed by \artifact{scripts/supplement\_cis.py} (re-run) into
\artifact{corpus/real\_world/supplement\_cis.json}: a nonparametric bootstrap
with $B=10{,}000$ resamples, seed $20260712$, resampling \emph{repositories}
(the \slug{owner\_\_repo} prefix of each slug) with replacement and taking all
of a drawn repository's workflows, which is the clustering the main text
promises. As a correctness check, the script's point estimates reproduce
\slug{collision\_free\_fidelity} exactly.

\begin{table}[htbp]
\centering\small
\caption{Corrected instrument (v2) on the matched, collision-free set,
with repository-clustered bootstrap 95\% CIs. $n$ = workflows, $r$ = distinct
repositories. Source: \artifact{corpus/real\_world/supplement\_cis.json}
(\slug{v2:*}); point estimates cross-checked against
\artifact{matched\_fidelity.json} $\to$
\slug{provenance\_collisions.collision\_free\_fidelity.v2\_corrected}.}
\label{tab:fidelity_v2}
\begin{tabular}{@{}lrrccc@{}}
\toprule
\textbf{Framework} & \textbf{$n$} & \textbf{$r$} &
\textbf{Edge P} & \textbf{Edge R} & \textbf{Kind acc.} \\
\midrule
LangGraph & 38 & 29 & 0.954 [0.865,\,1.000] & 0.836 [0.695,\,0.943] & 0.670 [0.606,\,0.738] \\
CrewAI    & 20 & 17 & 0.700 [0.440,\,0.944] & 0.692 [0.435,\,0.939] & 0.929 [0.850,\,1.000] \\
AutoGen   & 30 & 19 & 0.966 [0.917,\,1.000] & 0.778 [0.678,\,0.872] & 0.890 [0.832,\,0.945] \\
ADK       & 18 & 13 & 0.394 [0.108,\,0.678] & 0.343 [0.067,\,0.636] & 0.920 [0.854,\,0.979] \\
\midrule
\textbf{Overall} & \textbf{106} & \textbf{78} &
\textbf{0.814} [0.720,\,0.897] & \textbf{0.709} [0.613,\,0.798] & \textbf{0.824} [0.778,\,0.868] \\
\bottomrule
\end{tabular}
\end{table}

\begin{table}[htbp]
\centering\small
\caption{As-mined instrument (v1) on the same matched, collision-free set,
with repository-clustered bootstrap 95\% CIs. Only LangGraph differs from
Table~\ref{tab:fidelity_v2}; ADK, AutoGen, and CrewAI are bit-identical between
instruments. Source: \artifact{corpus/real\_world/supplement\_cis.json}
(\slug{v1:*}).}
\label{tab:fidelity_v1}
\begin{tabular}{@{}lrrccc@{}}
\toprule
\textbf{Framework} & \textbf{$n$} & \textbf{$r$} &
\textbf{Edge P} & \textbf{Edge R} & \textbf{Kind acc.} \\
\midrule
LangGraph & 38 & 29 & 0.954 [0.865,\,1.000] & 0.678 [0.549,\,0.791] & 0.659 [0.590,\,0.730] \\
CrewAI    & 20 & 17 & 0.700 [0.440,\,0.944] & 0.692 [0.435,\,0.939] & 0.929 [0.850,\,1.000] \\
AutoGen   & 30 & 19 & 0.966 [0.917,\,1.000] & 0.778 [0.678,\,0.872] & 0.890 [0.832,\,0.945] \\
ADK       & 18 & 13 & 0.394 [0.108,\,0.678] & 0.343 [0.067,\,0.636] & 0.920 [0.854,\,0.979] \\
\midrule
\textbf{Overall} & \textbf{106} & \textbf{78} &
\textbf{0.814} [0.720,\,0.897] & \textbf{0.652} [0.561,\,0.739] & \textbf{0.820} [0.773,\,0.866] \\
\bottomrule
\end{tabular}
\end{table}

\begin{table}[htbp]
\centering\small
\caption{Node detection on the matched, collision-free set with
repository-clustered bootstrap 95\% CIs (v2 corrected instrument). Node
detection is the stable dimension; the framework spread lives in edges and
kinds. Source: \artifact{corpus/real\_world/supplement\_cis.json}.}
\label{tab:nodes}
\begin{tabular}{@{}lrcc@{}}
\toprule
\textbf{Framework} & \textbf{$n$} & \textbf{Node P} & \textbf{Node R} \\
\midrule
LangGraph & 38 & 0.967 [0.903,\,1.000] & 0.930 [0.841,\,0.992] \\
CrewAI    & 20 & 0.867 [0.760,\,0.972] & 0.867 [0.760,\,0.972] \\
AutoGen   & 30 & 1.000 [1.000,\,1.000] & 0.955 [0.892,\,1.000] \\
ADK       & 18 & 0.713 [0.546,\,0.863] & 0.595 [0.400,\,0.782] \\
\midrule
\textbf{Overall} & \textbf{106} & \textbf{0.914} [0.866,\,0.956] & \textbf{0.868} [0.810,\,0.920] \\
\bottomrule
\end{tabular}
\end{table}

\paragraph{Reading the intervals.}
Three observations the main text has no room for. First, ADK's interval on edge
recall, $[0.067, 0.636]$, is enormous: it rests on 18 workflows in 13
repositories, and the point estimate should be treated as an order-of-magnitude
statement only. Second, CrewAI's edge precision and recall coincide exactly
($0.700$/$0.692$ with near-identical intervals) because the CrewAI extractor
emits a chain whose length is determined by the task list, so a missing task
costs symmetric precision and recall. Third, LangGraph's kind accuracy interval
$[0.606, 0.738]$ excludes every other framework's point estimate, which is the
statistical form of the paper's claim that LangGraph's residual error is a
\emph{kind} error rather than an \emph{edge} error.

\subsection{The unmatched breakdown}

The main table reports the matched, collision-free set ($n=106$) because it is
the only set on which the two instruments are scored over identical inputs. For
transparency we give the two larger, less defensible sets alongside it.

\begin{table}[htbp]
\centering\small
\caption{The three nested fidelity sets, as-mined instrument (v1). $n=119$ is
the original validation sample as reported before the keying defect was found
(120 minus the self-repo graph); $n=115$ is the v1$\cap$v2 intersection;
$n=106$ additionally removes the 9 provenance-colliding slugs. Sources:
\artifact{matched\_fidelity.json} keys
\slug{unmatched\_fidelity\_for\_reference.v1\_asmined\_n119\_style},
\slug{matched\_fidelity.v1\_asmined}, and
\slug{provenance\_collisions.collision\_free\_fidelity.v1\_asmined}.}
\label{tab:unmatched}
\begin{tabular}{@{}llrrrrr@{}}
\toprule
\textbf{Set} & \textbf{Framework} & \textbf{$n$} &
\textbf{Node P} & \textbf{Node R} & \textbf{Edge P} & \textbf{Edge R} \\
\midrule
Unmatched (119)
   & LangGraph & 40 & 0.968 & 0.910 & 0.956 & 0.682 \\
 & CrewAI    & 20 & 0.867 & 0.867 & 0.700 & 0.692 \\
 & AutoGen   & 35 & 1.000 & 0.942 & 0.957 & 0.770 \\
 & ADK       & 24 & 0.711 & 0.622 & 0.420 & 0.382 \\
 & \emph{Overall} & \emph{119} & \emph{0.909} & \emph{0.854} & \emph{0.805} & \emph{0.649} \\
\midrule
Matched (115)
   & LangGraph & 40 & 0.968 & 0.910 & 0.956 & 0.682 \\
 & CrewAI    & 20 & 0.867 & 0.867 & 0.700 & 0.692 \\
 & AutoGen   & 32 & 1.000 & 0.958 & 0.968 & 0.792 \\
 & ADK       & 23 & 0.699 & 0.606 & 0.395 & 0.355 \\
 & \emph{Overall} & \emph{115} & \emph{0.906} & \emph{0.855} & \emph{0.803} & \emph{0.649} \\
\midrule
\textbf{Collision-free (106)}
   & LangGraph & 38 & 0.967 & 0.912 & 0.954 & 0.678 \\
 & CrewAI    & 20 & 0.867 & 0.867 & 0.700 & 0.692 \\
 & AutoGen   & 30 & 1.000 & 0.955 & 0.966 & 0.778 \\
 & ADK       & 18 & 0.713 & 0.595 & 0.394 & 0.343 \\
 & \emph{Overall} & \emph{106} & \emph{0.914} & \emph{0.862} & \emph{0.814} & \emph{0.652} \\
\bottomrule
\end{tabular}
\end{table}

\paragraph{What the exclusions cost.}
The overall v1 numbers barely move across the three sets (edge recall $0.649
\to 0.649 \to 0.652$; node precision $0.909 \to 0.906 \to 0.914$), which is the
reassuring part. The unreassuring part is that ADK's edge recall falls
monotonically, $0.382 \to 0.355 \to 0.343$, while its $n$ falls $24 \to 23 \to
18$. The exclusions are not uniform, and Section~\ref{sec:collisions}
quantifies the bias.

\paragraph{Current CI authority.}
\artifact{corpus/real\_world/confidence\_intervals.json} and
\artifact{supplement\_cis.json} now both use the matched, collision-free
$n=106$ slug set and repository-clustered bootstrap. The first additionally
records the self-repository-excluded 186-flag triage and HGP-1 intervals. Their
point estimates are checked against
\artifact{matched\_fidelity.json}, and
\artifact{validate\_submission\_claims.py} fails if the old unmatched
$n=120/187$ inputs reappear.

% ===========================================================================
\section{Full Triage Breakdown for the 186 Flags}
\label{sec:triage}

\emph{Promised at} \artifact{sections/03\_realworld.tex}, paragraph
\emph{Flag precision on extracted graphs}: ``(full triage breakdown in the
supplement)''.

\subsection{Provenance and the 186-vs-187 discrepancy}

The triage universe lives at
\artifact{corpus/real\_world/validated\_results.json}, key
\slug{triage\_details}, which holds \textbf{187} records. Exactly one of them
has a slug naming our own repository (redacted here for anonymity), suffix
\slug{\_\_test\_extract\_adk}---a graph from
our own test fixtures, the contamination the main text describes as detected
during revision and excluded from every number. Removing it leaves
\textbf{186}, and the label totals then reproduce the main text exactly. All
tables in this section are computed on the 186; the committed
\artifact{confidence\_intervals.json} was computed \emph{before} that exclusion
and reports the 187-flag figures ($42.8\%$/$49.7\%$ instead of
$42.5\%$/$50.0\%$), which is why the main text does not cite it here and why
those numbers must not be lifted from earlier drafts.

\begin{table}[htbp]
\centering\small
\caption{Per-check $\times$ per-label counts for all 186 triaged flags.
Rows are the checks as recorded in the artifact (two rows carry the legacy
aliases that \artifact{error\_decomposition.json} maps to
\slug{router\_shape} and \slug{human\_gate\_coverage} respectively).
Source: \artifact{corpus/real\_world/validated\_results.json} $\to$
\slug{triage\_details}, self-repo record excluded.}
\label{tab:triage}
\begin{tabular}{@{}lrrrrr@{}}
\toprule
\textbf{Check} & \textbf{extraction} & \textbf{intentional} & \textbf{arguable} &
\textbf{real} & \textbf{Total} \\
 & \textbf{artifact} & & & \textbf{defect} & \\
\midrule
\slug{human\_presence}          & 8  & 93 & 10 & 2 & 113 \\
\slug{reverse\_reachability}    & 25 & 0  & 0  & 0 & 25 \\
\slug{dead\_ends}               & 24 & 0  & 1  & 0 & 25 \\
\slug{exit\_reachability}       & 14 & 0  & 0  & 0 & 14 \\
\slug{tool\_declarations}       & 5  & 0  & 0  & 0 & 5 \\
\slug{router\_shape}            & 2  & 0  & 0  & 0 & 2 \\
\slug{router-non-conditional-edges}$^{\ast}$ & 1 & 0 & 0 & 0 & 1 \\
\slug{sensitive-path-bypasses-human-review}$^{\dagger}$ & 0 & 0 & 1 & 0 & 1 \\
\midrule
\textbf{Total} & \textbf{79} & \textbf{93} & \textbf{12} & \textbf{2} & \textbf{186} \\
\textbf{Share} & 42.5\% & 50.0\% & 6.5\% & 1.1\% & 100\% \\
\bottomrule
\end{tabular}\\[3pt]
{\footnotesize $^{\ast}$legacy alias for \slug{router\_shape} (combined: 3
flags). $^{\dagger}$legacy alias for \slug{human\_gate\_coverage}.}
\end{table}

\paragraph{Consistency with the main text.}
Table~\ref{tab:triage} reproduces every headline of
\artifact{sections/03\_realworld.tex}: 186 flags; label shares
$42.5\%$/$50.0\%$/$6.5\%$/$1.1\%$; 2 classified as genuine, both under
\slug{human\_presence} on sensitive workflows; 12 \textsc{arguable}; and
$113/186=61\%$ of flags from the blunt \slug{human\_presence} check alone. The
72 structural flags are the sum of the six structural rows
($25+25+14+5+2+1=72$), of which 71 are extraction artifacts and 1 is arguable,
with \emph{zero} classified as genuine---the main text's ``none of the 72
structural flags was classified as genuine''.

\subsection{The four-way attribution, by check}

Attribution maps each triage label to a cause under the rules recorded in
\artifact{corpus/real\_world/error\_decomposition.json} $\to$
\slug{\_meta.attribution\_rules}: \slug{extraction\_artifact} and
\slug{gt\_error} $\to$ \textsc{extractor}; \slug{intentional} $\to$
\textsc{policy-spec}; \slug{arguable} $\to$ \textsc{label};
\slug{real\_defect} $\to$ \textsc{actionable}. (No record in the 186 carries
the \slug{gt\_error} label, so \textsc{extractor} coincides exactly with
\slug{extraction\_artifact} here.) \textsc{checker} is not separable in this
\emph{label-derived} attribution and is folded into \textsc{policy-spec}; both
are on the check side, so the extraction-versus-check comparison is unaffected.
Section~\ref{sec:survival} separates the two with a second, label-independent
instrument and finds \textsc{checker} $= 1.1\%$ $[0.3,3.8]$ --- i.e.\ the
\textsc{policy-spec} bucket reported here is $\approx\!98\%$ specification
error, and the two instruments agree on $93.0\%$ of flags.

\begin{table}[htbp]
\centering\small
\caption{Four-way attribution by check, with Wilson 95\% CIs (\%). Source:
\artifact{corpus/real\_world/error\_decomposition.json} $\to$
\slug{decomposition\_186.by\_check}. Checks with a single flag are reported
for completeness; their intervals are uninformative.}
\label{tab:attr_check}
\begin{tabular}{@{}lrllll@{}}
\toprule
\textbf{Check} & \textbf{$n$} & \textbf{Extractor} & \textbf{Policy-spec} &
\textbf{Label} & \textbf{Actionable} \\
\midrule
\slug{human\_presence} & 113
 & 7.1 [3.6,\,13.4] & 82.3 [74.2,\,88.2] & 8.8 [4.9,\,15.5] & 1.8 [0.5,\,6.2] \\
\slug{reverse\_reachability} & 25
 & 100 [86.7,\,100] & --- & --- & --- \\
\slug{dead\_ends} & 25
 & 96.0 [80.5,\,99.3] & --- & 4.0 [0.7,\,19.5] & --- \\
\slug{exit\_reachability} & 14
 & 100 [78.5,\,100] & --- & --- & --- \\
\slug{tool\_declarations} & 5
 & 100 [56.5,\,100] & --- & --- & --- \\
\slug{router\_shape} & 3
 & 100 [43.9,\,100] & --- & --- & --- \\
\slug{human\_gate\_coverage} & 1
 & --- & --- & 100 [20.6,\,100] & --- \\
\midrule
\textbf{All} & \textbf{186}
 & \textbf{42.5} [35.6,\,49.7] & \textbf{50.0} [42.9,\,57.1]
 & \textbf{6.5} [3.7,\,10.9] & \textbf{1.1} [0.3,\,3.8] \\
\bottomrule
\end{tabular}
\end{table}

\subsection{The four-way attribution, by framework}

\begin{table}[htbp]
\centering\small
\caption{Four-way attribution by framework, with Wilson 95\% CIs (\%). Source:
\artifact{corpus/real\_world/error\_decomposition.json} $\to$
\slug{decomposition\_186.by\_framework}.}
\label{tab:attr_fw}
\begin{tabular}{@{}lrllll@{}}
\toprule
\textbf{Framework} & \textbf{$n$} & \textbf{Extractor} & \textbf{Policy-spec} &
\textbf{Label} & \textbf{Actionable} \\
\midrule
LangGraph & 111 & \textbf{65.8} [56.5,\,73.9] & 28.8 [21.2,\,37.9] & 5.4 [2.5,\,11.3] & --- \\
AutoGen   & 31  & 16.1 [7.1,\,32.6]  & \textbf{74.2} [56.8,\,86.3] & 9.7 [3.4,\,24.9] & --- \\
ADK       & 24  & 4.2 [0.7,\,20.2]   & \textbf{83.3} [64.1,\,93.3] & 8.3 [2.3,\,25.9] & 4.2 [0.7,\,20.2] \\
CrewAI    & 20  & ---                & \textbf{90.0} [69.9,\,97.2] & 5.0 [0.9,\,23.6] & 5.0 [0.9,\,23.6] \\
\midrule
\textbf{All} & \textbf{186} & 42.5 [35.6,\,49.7] & 50.0 [42.9,\,57.1] & 6.5 [3.7,\,10.9] & 1.1 [0.3,\,3.8] \\
\bottomrule
\end{tabular}
\end{table}

\paragraph{The framework split is the aggregate's whole story.}
LangGraph is the only framework in which extraction is the majority cause
($65.8\%$), and it contributes $111/186 = 60\%$ of all flags. In the other three
frameworks policy misspecification dominates by a wide margin ($74$--$90\%$),
and extraction attribution is at or near zero. This is a direct consequence of
the structural-check asymmetry the main text notes: the CrewAI, AutoGen, and ADK
extractors emit connected topologies by construction and \emph{cannot} produce a
dead end or an unreachable exit, so essentially all of their flags come from the
blunt human-presence check. The aggregate $42.5\%$ vs.\ $50.0\%$ is therefore not
a statement about extraction in general; it is a weighted average of one
extraction-dominated framework and three policy-dominated ones.

\subsection{By check family, and the symmetric ablation}

\begin{table}[htbp]
\centering\small
\caption{Extraction versus check-side attribution among the 184 non-actionable
flags, split by check family, with Wilson 95\% CIs (\%). Source:
\artifact{error\_decomposition.json} $\to$
\slug{decomposition\_186.by\_check\_family}.}
\label{tab:family}
\begin{tabular}{@{}lrll@{}}
\toprule
\textbf{Family} & \textbf{$n$ (non-act.)} & \textbf{Extraction-attributable} &
\textbf{Check-side} \\
\midrule
Structural (6 checks) & 72  & \textbf{98.6\%} [92.5,\,99.8] & 0.0\% [0.0,\,5.1] \\
Human-gate            & 112 & 7.1\% [3.7,\,13.5]            & \textbf{83.0\%} [75.0,\,88.9] \\
\bottomrule
\end{tabular}
\end{table}

\begin{table}[htbp]
\centering\small
\caption{The symmetric ablation underlying the main text's ``what survives a
perfect graph?'' paragraph. Source: \artifact{error\_decomposition.json} $\to$
\slug{symmetric\_ablation} and \slug{headline}.}
\label{tab:ablation}
\begin{tabular}{@{}lrr@{}}
\toprule
\textbf{Quantity} & \textbf{v1 extraction} & \textbf{Reference (oracle) graph} \\
\midrule
Workflows                                & 119 & 119 \\
Total flags                              & 184 & 122 \\
Flags per workflow                       & 1.55 & 1.03 \\
Structural flags                         & 71 & 13 \\
Workflows with $\ge 1$ structural flag    & 26 & 11 \\
\slug{human\_presence} firings         & 113 & 106 \\
\slug{human\_gate\_coverage} firings   & 0 & 3 \\
\bottomrule
\end{tabular}\\[3pt]
{\footnotesize Flags removed by a perfect graph: 62 ($33.7\%$ of flag volume).
Residual that a perfect graph still produces: 122 ($66.3\%$). Genuine defects
recovered only on the reference graphs: 2.}
\end{table}

Two details worth stating. First, the ablation's baseline is \textbf{184} flags,
not 186, because it re-runs the check suite on the 119 workflows rather than
replaying the triage records; the two universes differ by two flags and the main
text is careful never to pool them. Second, \slug{human\_gate\_coverage} fires
\emph{zero} times on extracted graphs and \emph{three} times on reference
graphs---the risk-aware check is not insensitive, it is starved of declared
bindings, which is the paper's central diagnosis expressed as a single row.

% ===========================================================================
\section{The Survival Test: Separating \textsc{checker} from \textsc{policy-spec},
and Decomposing Misses}
\label{sec:survival}

The four-way attribution above is derived from triage \emph{labels}, and those
labels cannot isolate \textsc{checker}: the triager was asked ``is this a real
defect?'', not ``whose fault is this flag?''. This section reports a second,
\emph{label-independent} instrument that closes that gap and extends the same
logic to false negatives. It is implemented in
\artifact{scripts/fp\_fn\_decomposition.py} and emits
\artifact{corpus/real\_world/fp\_fn\_decomposition.json}. It re-reads no source
code: every input is a committed artifact, so the decomposition replays
deterministically.

\subsection{The rule, stated before it was run}

For each triaged flag (workflow $s$, check $C$) judged non-actionable:

\begin{description}
\item[R1] Re-run $C$ on the source-reconstructed reference graph
  $\mathrm{GT}(s)$. If $C$ does \emph{not} fire, the reconstruction removed the
  flag $\Rightarrow$
  \textsc{extractor}.
\item[R2] If $C$ fires and $C$ is \emph{structural}
  (\slug{dead\_ends}, \slug{exit\_reachability}, \slug{reverse\_reachability},
  \slug{router\_shape}, \slug{tool\_declarations}) $\Rightarrow$
  \textsc{checker}. A structural check is a claim about topology alone; it has
  no per-workflow applicability parameter to misfire on. If the topology is
  faithful and the finding is still not a defect, the decision procedure is
  wrong for that framework's semantics.
\item[R3] If $C$ fires and $C$ is a \emph{policy} check
  (\slug{human\_presence}, \slug{human\_gate\_coverage}) $\Rightarrow$
  \textsc{policy-spec}, except when $\mathrm{GT}(s)$ already contains a node of
  kind \slug{human}, where the check contradicts evidence present in the
  faithful graph $\Rightarrow$ \textsc{checker}.
\item[R4/R5] \slug{arguable} $\Rightarrow$ \textsc{label};
  \slug{real\_defect} $\Rightarrow$ actionable (outside the denominator).
  These precede R1--R3: a disputed label is a label error whatever the graph
  says.
\end{description}

\paragraph{The validity guard.} R1--R3 presume the reference graph is correct
for the predicate under test. Our own results say it sometimes is not, and two
committed sources record exactly where: (a)~the GT triage pass labelled the flag
\slug{gt\_error}, i.e.\ the reconstruction misrepresents the source; (b)~the
verdict turns on a \slug{human} node that the HGP-1 placebo audit found
non-discharging. When the guard trips, the survival test is not evidence about
the check, and we record \textsc{reference-error} rather than force the flag
onto a side. Reporting it as its own cause is the point: it measures how often
the source-reconstructed reference is itself wrong, which on this sample is
\textbf{5 of 186} flags ($2.7\%$). Suppressing that category would have
manufactured 3 spurious \textsc{checker} findings and 2 spurious
\textsc{extractor} findings.

\subsection{False positives: the checker is not the problem}

\begin{table}[htbp]
\centering\small
\caption{Label-independent decomposition of the 186 triaged flags by the
survival test, with Wilson 95\% CIs (\%). Source:
\artifact{fp\_fn\_decomposition.json} $\to$
\slug{false\_positives.over\_all\_flags}.}
\label{tab:survival}
\begin{tabular}{@{}lrrl@{}}
\toprule
\textbf{Cause} & \textbf{$n$} & \textbf{\%} & \textbf{Wilson 95\%} \\
\midrule
\textsc{policy-spec}      & 94 & 50.5 & [43.4,\,57.6] \\
\textsc{extractor}        & 71 & 38.2 & [31.5,\,45.3] \\
\textsc{label}            & 12 &  6.5 & [3.7,\,10.9]  \\
\textsc{reference-error}  &  5 &  2.7 & [1.2,\,6.1]   \\
\textsc{checker}          &  2 &  1.1 & [0.3,\,3.8]   \\
Actionable                &  2 &  1.1 & [0.3,\,3.8]   \\
\bottomrule
\end{tabular}
\end{table}

\begin{table}[htbp]
\centering\small
\caption{The same decomposition split by check family, over the 184
non-actionable flags. Source: \slug{false\_positives.by\_check\_family}.}
\label{tab:survival_family}
\begin{tabular}{@{}lrllll@{}}
\toprule
\textbf{Family} & \textbf{$n$} & \textbf{Extr.} & \textbf{Checker} &
\textbf{Policy} & \textbf{Ref./Label} \\
\midrule
Structural & 72  & \textbf{94.4\%} & 2.8\% & 0.0\% & 2.8\% \\
Human-gate & 112 & 2.7\%           & 0.0\% & \textbf{83.9\%} & 13.4\% \\
\bottomrule
\end{tabular}
\end{table}

Three things follow. First, \textsc{checker} is now \emph{quantified} rather
than merely bounded below: $2/186 = 1.1\%$ $[0.3,3.8]$. The two instances are
the LangGraph \slug{interrupt()} multi-turn pattern, where a genuinely
outgoing-edge-free node is a deliberate re-entry point rather than a dead end,
and a LangGraph \slug{ToolNode} whose tools are fetched from an MCP server at
runtime and so are statically undeclarable. Both are framework-semantics gaps in
the check, and neither is a policy question.

Second, the check-side error the paper reports is therefore almost entirely
\emph{specification}, not \emph{logic}: $50.5\%$ \textsc{policy-spec} against
$1.1\%$ \textsc{checker}. The checks do what they say; what they say is wrong
for these workflows. This matters for the paper's prescription---writing a
better checker is not the lever; instantiating the policy per workflow is.

Third, the two instruments agree. Mapping the committed triage labels
directly (\slug{extraction\_artifact}$\to$\textsc{extractor},
\slug{intentional}$\to$\textsc{policy-spec}) reproduces the survival test's
verdict on \textbf{173 of 186} flags, $93.0\%$. The 13 disagreements are listed
in the artifact with their triggering rule; 5 are guard trips, and the
remainder are cases where a \slug{human} node present in the reconstruction but
absent from the extraction (or vice versa) moves the flag between
\textsc{extractor} and \textsc{policy-spec}. This is not independent human
validation and is not offered as a substitute for it. It is a weaker but real
claim: the headline decomposition does not depend on the LLM triage labels,
because a mechanical procedure over graphs alone reproduces it.

\subsection{False negatives: the opposite budget}

Flag triage is blind to misses by construction. The survival test extends to
them directly: hold the policy fixed, vary the graph, and ask which of the 12
confirmed HGP-1 violations the pipeline reports.

\begin{table}[htbp]
\centering\small
\caption{Why the pipeline misses each of the 12 confirmed human-gate
violations, under the as-mined extractor and under the reference graph. Source:
\slug{false\_negatives}.}
\label{tab:fn}
\begin{tabular}{@{}lrr@{}}
\toprule
\textbf{Cause of the miss} & \textbf{As-mined} & \textbf{Reference graph} \\
\midrule
Detected (no miss)                & 0  & 1 \\
\textsc{extractor}                & \textbf{12} & --- \\
\textsc{abstraction}              & --- & \textbf{9} \\
\textsc{reference-error}          & --- & 1 \\
\textsc{checker}                  & 0  & 1 \\
\midrule
\textbf{Total missed}             & \textbf{12} & \textbf{11} \\
\bottomrule
\end{tabular}
\end{table}

The as-mined column is unmixed: \emph{no} mined graph declares a sensitive tool
at all, so the risk-aware check could not have fired on any of the 12, and every
miss is extraction. The reference-graph column is the informative one. Swapping
in the source-reconstructed reference recovers \textbf{1 of 12}. Of the 11 that
remain:

\begin{itemize}
\item \textbf{9 are \textsc{abstraction}}: the effect is absent from the
  \emph{reconstruction} too, because it lives inside a node \emph{body}. No
  extractor could recover it within the \slug{NodeKind}/\slug{EdgeKind}/declared-tool
  vocabulary these graphs use; the fix is a richer abstraction, not a better
  reader of the same one.
\item \textbf{1 is \textsc{reference-error}}: the reconstruction omitted tools
  the source declares statically (the \slug{Harsh1210\_\_llm-opsbot\_\_main}
  \slug{gt\_error}), so this one \emph{is} recoverable by a better model.
\item \textbf{1 is \textsc{checker}}: the workflow was cleared on a gate that
  cannot block---the placebo failure mode, appearing here as a false negative
  rather than a false positive.
\end{itemize}

\paragraph{Why this is the paper's sharpest result.} The two directions have
\emph{opposite} error budgets. Non-actionable findings are check-side
($50.5\%$ policy specification versus $38.2\%$ extraction). Misses are
model-side ($100\%$ extraction as mined). But the reconstruction repairs only
$1/12$ of the misses, because the missing information sits \emph{below} the
abstraction rather than above it. A reader who concluded from the false-positive
decomposition that ``extraction fidelity is a secondary concern'' would draw
exactly the wrong lesson: extraction fidelity bounds what the analysis can
\emph{see}, and improving the extractor within a body-blind abstraction moves
that bound by one workflow in twelve.

This is the same boundary the certification gate meets
(Section~\ref{sec:h1}). Exact provenance certifies that an \emph{edge} was read
correctly; it says nothing about the events inside the node that edge connects.
Table~\ref{tab:fn} is the empirical measurement of that gap: on 9 of 12
source-audited violations, the abstraction is silent about the very effect the policy
regulates.

% ===========================================================================
\section{Scalability, Extractor Accuracy, and a Modeling-Effort Comparison}
\label{sec:eval}

\emph{Promised at} \artifact{sections/04\_results.tex}: ``details, per-framework
extractor accuracy on stubs, and a Promela/SMV modeling-effort comparison are in
the supplementary material.''

\subsection{Scalability (re-run)}

\begin{table}[htbp]
\centering\small
\caption{Structural check time versus graph size, \textbf{freshly re-run} for
this supplement via \artifact{scripts/benchmark\_scale.py} (median of 10 trials,
synthetic graphs at edge density 2.0). Output artifact:
\artifact{scripts/scaling\_results.json}. Monitor compilation and evaluation are
size-independent and are reported once below the table.}
\label{tab:scaling}
\begin{tabular}{@{}rrrr@{}}
\toprule
\textbf{Requested nodes} & \textbf{Actual nodes} & \textbf{Actual edges} &
\textbf{Structural checks (ms)} \\
\midrule
50    & 52    & 104     & 0.060 \\
100   & 102   & 214     & 0.111 \\
200   & 202   & 430     & 0.221 \\
500   & 502   & 1{,}074  & 0.592 \\
1{,}000 & 1{,}002 & 2{,}128  & 1.267 \\
2{,}000 & 2{,}002 & 4{,}252  & 2.977 \\
5{,}000 & 5{,}002 & 10{,}713 & \textbf{7.191} \\
\bottomrule
\end{tabular}\\[3pt]
{\footnotesize Monitor compilation (5 rules): $0.259$\,ms. Monitor evaluation
(1{,}000 events $\times$ 5 rules): $9.264$\,ms, i.e.\ ${\approx}108{,}000$
events/s. Both are independent of graph size.}
\end{table}

\paragraph{These numbers are hardware-dependent, and we say so loudly.}
The same script, same code, and same machine produced $7.611$\,ms at 5{,}000
nodes on one invocation and $7.191$\,ms on the next---a $6\%$ run-to-run spread
from nothing but scheduling noise. More importantly, the previously committed
artifact recorded $239.838$\,ms and an earlier arXiv draft reported
$104.7$\,ms at the same size. The three figures span a factor of ${\approx}33$,
and none of them is wrong: they are three different machines. \textbf{Only the
asymptotic shape should be read as a property of the system}; the absolute
milliseconds are a property of whatever hardware ran the benchmark, and we
therefore make no cross-draft comparison of them.

\paragraph{Shape.} Over $50\to5{,}000$ nodes the graph grows $96\times$ in
$|V|+|E|$ while check time grows $120\times$, i.e.\ very slightly super-linear.
The excess comes from quadratic edge scanning in the router-shape and
tool-declaration checks, which is invisible at the sizes that actually occur: the
mined corpus has a \emph{median of 2 non-sentinel nodes}
(\artifact{sections/03\_realworld.tex}; \artifact{corpus/real\_world/metadata\_v2.json}),
so the largest row of Table~\ref{tab:scaling} is roughly three orders of
magnitude beyond anything we observed in the wild.

\subsection{Per-framework extractor accuracy}
\label{sec:stubs}

\paragraph{Disclosure: the promised ``on stubs'' measurement is not available.}
The main text promises ``per-framework extractor accuracy on stubs''. The
harness for it is \artifact{scripts/extractor\_accuracy.py}, whose defaults point
at \artifact{corpus/graphs/} (extracted) and \artifact{corpus/ground\_truth/}
(reference). \artifact{corpus/ground\_truth/} \textbf{does not exist in the
repository}, and the script exits with ``No ground truth files found''. The
nearest substitute, running \artifact{corpus/graphs/} against
\artifact{corpus/curated/}, returns $1.000$ on every metric for all four
frameworks---but that is an artifact, not a result: \artifact{corpus/curated/} is
\emph{authored} by \artifact{scripts/build\_corpus.py} as literal
\slug{GraphNode}/\slug{GraphEdge} constructions, and \artifact{corpus/graphs/}
is a copy of it differing only in the provenance fields. Comparing them measures
a file copy. \textbf{We therefore report no stub-accuracy table and record the
promise as unmet rather than filling it with a circular $1.000$.} Producing it
honestly requires stub \emph{source files} plus independent reference
annotations; neither is committed.

\paragraph{What we can report: accuracy on real mined code.}
The same script, run against the real-world corpus (re-run:
\slug{--corpus corpus/real\_world/graphs --truth corpus/real\_world/ground\_truth},
output \artifact{scripts/accuracy\_results.json}), gives a genuine per-framework
measurement. Note this is the \emph{unmatched} $n=120$ set including the
self-repo graph, so it is a per-instrument diagnostic and not the paper's
headline; Section~\ref{sec:fidelity} carries the defensible numbers.

\begin{table}[htbp]
\centering\small
\caption{Per-framework extractor accuracy against the LLM-reconstructed
reference graphs (re-run; \artifact{scripts/accuracy\_results.json}). Unmatched
$n=120$ set. ``P''/``R'' are macro-averages over workflows.}
\label{tab:extacc}
\begin{tabular}{@{}lrrrrrr@{}}
\toprule
\textbf{Framework} & \textbf{$n$} & \textbf{Node P} & \textbf{Node R} &
\textbf{Kind acc.} & \textbf{Edge P} & \textbf{Edge R} \\
\midrule
LangGraph & 40 & 0.968 & 0.910 & 0.647 & 0.956 & 0.682 \\
CrewAI    & 20 & 0.867 & 0.867 & 0.929 & 0.700 & 0.692 \\
AutoGen   & 35 & 1.000 & 0.942 & 0.900 & 0.957 & 0.770 \\
ADK       & 25 & 0.709 & 0.602 & 0.943 & 0.403 & 0.367 \\
\midrule
\textbf{Overall} & \textbf{120} & \textbf{0.907} & \textbf{0.848} &
\textbf{0.829} & \textbf{0.798} & \textbf{0.644} \\
\bottomrule
\end{tabular}
\end{table}

\paragraph{Failure modes.}
The same run categorises every node-kind mismatch. Over all 120 workflows there
are \textbf{104} categorised mismatches, distributed as
\slug{custom\_node\_type} \textbf{98} (94.2\%) --- ``unrecognized node kinds or
custom agent types'', which in practice is overwhelmingly the
\slug{tool}$\rightarrow$\slug{llm} error --- and
\slug{naming\_heuristic\_failure} \textbf{6} (5.8\%) --- ``human/router
detection by name convention failed''. Two further categories exist in the
taxonomy (\slug{dynamic\_construction}, \slug{version\_incompatibility}) and
score \textbf{0}. This single distribution is the quantitative core of the
paper's diagnosis: kind error is not a long tail of framework quirks, it is one
failure mode with a 94\% share, and it is the one that hides side-effecting
tools from the risk-aware gate (Section~\ref{sec:h1}).

\subsection{Modeling-effort comparison against general-purpose model checkers}

The table below is reproduced verbatim from the committed artifact
\artifact{papers/paper1/generated/comparison\_table.tex} (generated by
\artifact{scripts/generate\_comparison\_table.py}). It is qualitative by
construction and we do not dress it up as a measurement.

\begin{table}[htbp]
\centering\small
\caption{Comparison with general-purpose model checkers. Source:
\artifact{papers/paper1/generated/comparison\_table.tex}.}
\label{tab:comparison_mc}
\resizebox{\textwidth}{!}{%
\begin{tabular}{@{}llllll@{}}
\toprule
\textbf{Tool} & \textbf{Input} & \textbf{Properties} & \textbf{Modeling} &
\textbf{Time} & \textbf{Domain} \\
\midrule
SPIN~\citep{holzmann1997spin}   & Promela (manual) & Full LTL & High & Fast (small) & None \\
NuSMV~\citep{cimatti2002nusmv}  & SMV (manual)     & CTL + LTL & High & Fast (small) & None \\
CBMC~\citep{cbmc}               & C source         & Assertions & Medium & Bounded & None \\
\textbf{This work}              & Auto-extracted   & Safety LTL fragment &
\textbf{None} & $O(|V|{+}|E|)$ & \textbf{Agent workflows} \\
\bottomrule
\end{tabular}}\\[3pt]
{\footnotesize SPIN, NuSMV, and CBMC require manual translation of agent
workflows into their input languages. Our analyzer extracts models automatically
from framework APIs, eliminating the modeling step entirely.}
\end{table}

\paragraph{What ``Modeling: None'' does and does not mean, and what is not measured.}
Two honest qualifications. First, the modeling effort is not eliminated, it is
\emph{amortised}: the four per-framework extractors encode framework semantics
and were hand-written once, and Section~\ref{sec:fidelity} is the measurement of
how lossily they do it. What is eliminated is the \emph{per-workflow} manual
translation SPIN and NuSMV require. Second, the ``Time'' column is not a
head-to-head benchmark. The repository contains exactly one hand-written Promela
model, \artifact{corpus/comparisons/email\_triage.pml}, and \emph{no committed
SPIN or NuSMV run, no timing, and no lines-of-model count for any other
workflow}. \textbf{A quantitative modeling-effort comparison---person-minutes or
lines of Promela per workflow, across a sample---was not performed, and no
number for it appears in this supplement.} The comparison above is a
capability-and-workflow comparison only.

% ===========================================================================
\section{Sampling Design and the Limits of the Corpus}
\label{sec:sampling}

Not promised by the main text, but included because every prevalence number in
Section~3 of the paper is bounded by it. The following inlines
\artifact{papers/paper1/aaai/SAMPLING\_DESIGN.md}, which was reconstructed
entirely from repository artifacts.

\paragraph{Bottom line.}
The 922-graph corpus is a \textbf{search-ranked convenience corpus}, not a
probability sample of GitHub agent workflows. The 119-graph validation set is a
\textbf{hand-fixed, unseeded, quota-shaped subset} of it, hard-coded into a
driver script; no sampler, no seed, and no reproducible selection rule exists in
the repository. Inclusion probabilities are therefore not merely unestimated,
they are \textbf{undefined}. Every framework-share reweighting in the paper is a
\textbf{weighted descriptive statistic over a convenience corpus}, not a
design-based prevalence estimate for GitHub.

\subsection{The frame, verbatim}

The frame is the set of files returned by 11 GitHub code-search queries,
truncated by a per-query result cap and a global repository-list cap, in
GitHub's opaque relevance (``best match'') order, then all \slug{*.py} files
inside the repositories that survived that truncation. Queries are from
\artifact{scripts/mine\_github\_gh.py}:37--49, executed via
\slug{gh search code}.

\begin{table}[htbp]
\centering\small
\caption{The 11 verbatim mining queries. Source:
\artifact{scripts/mine\_github\_gh.py}:37--49.}
\label{tab:queries}
\begin{tabular}{@{}lp{0.62\linewidth}@{}}
\toprule
\textbf{Framework} & \textbf{Query string (verbatim)} \\
\midrule
LangGraph & \slug{StateGraph add\_node language:python} \\
LangGraph & \slug{StateGraph add\_conditional\_edges language:python} \\
LangGraph & \slug{from langgraph.graph add\_edge language:python} \\
CrewAI    & \slug{from crewai Crew Task language:python} \\
CrewAI    & \slug{crewai "process=Process" language:python} \\
AutoGen   & \slug{GroupChat autogen language:python} \\
AutoGen   & \slug{RoundRobinGroupChat language:python} \\
AutoGen   & \slug{SelectorGroupChat language:python} \\
ADK       & \slug{SequentialAgent google.adk language:python} \\
ADK       & \slug{from google.adk Agent language:python} \\
ADK       & \slug{ParallelAgent google.adk language:python} \\
\bottomrule
\end{tabular}
\end{table}

\paragraph{Caps and truncation order.}
\begin{itemize}
\item \textbf{Per-query result cap: 40.} \slug{gh search code ... --limit 40}
(\artifact{mine\_github\_gh.py}:52--58; default \slug{--per-query 40} at
line 295). Confirmed by artifact: \artifact{candidates.json} holds \textbf{341}
items over 11 queries, and AutoGen alone contributes exactly $120 = 3\times40$.
\item \textbf{No sort key is ever passed}, so \slug{gh search code} returns
GitHub's default ``best match'' relevance ranking---a proprietary, undocumented,
non-stationary function. Every subsequent truncation inherits this ordering.
\item \textbf{Repository-list truncation: \slug{unique[:max\_repos]}},
default 120, run at 180, at \artifact{mine\_github\_gh.py}:125 and :231. The
unique-repo list is built by first appearance in \artifact{candidates.json}
order, i.e.\ in relevance order, then cut. Verified: \artifact{repos.lg\_crew.json}
(180 repos) is set-identical to the first 180 unique repositories of
\artifact{candidates.json} in candidate order.
\item \textbf{No per-repository file cap.} \slug{find\_workflow\_files}
(\artifact{scripts/scrape\_workflows.py}:139--151) walks every \slug{*.py} file
under the clone, skipping only dot-dirs, \slug{venv}/\slug{.venv},
\slug{node\_modules}, \slug{\_\_pycache\_\_}. The paper's earlier phrase
``per-repository caps'' was inaccurate: the cap is on the repository \emph{list},
not on files per repository.
\item \textbf{Framework assignment is a text heuristic}, not a dependency check:
a file is labeled with framework $f$ if $\ge 2$ of $f$'s marker tokens appear
anywhere in it (\artifact{scrape\_workflows.py}:126--136).
\end{itemize}

\paragraph{Two mining runs.} Run~1 (LangGraph/CrewAI) used the full 341-candidate
list with \slug{--max-repos 180}; AutoGen and ADK extraction did not yet exist,
so the 180-repo prefix yielded 282 records over 149 repos. Run~2 applied the
\slug{--frameworks autogen,adk} filter (\artifact{mine\_github\_gh.py}:326--330)
over 139 unique repos---the 180 cap was not binding---for 798 records. Union:
318 candidate repositories, of which 253 yielded $\ge1$ extracted graph. Mining
date: \textbf{11 July 2026}, established independently in
\artifact{corpus/real\_world/MINING\_DATE\_FORENSICS.md}.

\paragraph{An undocumented silent de-duplication.}
Graph filenames are $\text{\slug{slug}} = \langle\text{owner}\rangle\_\_%
\langle\text{repo}\rangle\_\_\langle\text{file stem}\rangle$
(\artifact{mine\_github\_gh.py}:170, :256) written with a plain overwrite. Two
files with the same basename in one repository collapse to one graph,
last-writer-wins. Extracted records across both metadata files: \textbf{1080};
distinct slugs on disk: \textbf{930}; hence \textbf{150 extracted graphs
(13.9\%) were silently overwritten}. Worst case: \slug{ed-donner/agents}
produced \textbf{55} files that all collapsed into the single slug
\slug{ed-donner\_\_agents\_\_sidekick}. Section~\ref{sec:collisions} treats the
downstream consequence for the fidelity comparison.

\subsection{The 119-graph validation sample: no seed, no sampler}

The validation sample is not computed anywhere. It exists only as a hard-coded
JavaScript array literal: \artifact{scripts/wf\_run.js}, \slug{const GT}
(60 items: LangGraph 40, CrewAI 20) and \artifact{scripts/wf\_run\_v2.js}
(60 items: AutoGen 35, ADK 25). Total 120; the ADK self-repo item is dropped
downstream, giving 119 (LangGraph 40, CrewAI 20, AutoGen 35, ADK 24) from 87
repositories. \textbf{There is no sampler script, no manifest, and no record of
how the 120 were chosen.} The absolute paths embedded in those arrays point to a
different machine and a differently named repository, so the list was produced
outside this repository's history.

\paragraph{Refuted selection hypotheses.} Tested per framework against the corpus
in mining order, all \emph{negative}: first-$k$ prefix in mining order; first-$k$
after dropping $<3$-node graphs; first-$k$ with one graph per repository;
top-$k$ by graph size; round-robin over repositories in mining order (best
overlap $15/25$, ADK). Selected positions are scattered across each framework's
whole pool (LangGraph ranks 2--391 of 402; ADK 0--52 of 53), and the sample
contains up to 4 graphs from one repository, which rules out one-per-repo
designs.

\paragraph{No historical seed exists.}
\artifact{scripts/make\_annotation\_samples.py} (\slug{--seed}, default
\slug{20260713}) draws the \emph{human-annotation} sub-samples \emph{from} the
119, inheriting them as a fixed given, and cannot document how the 119 arose.
This remains true of every numeric result reported from that set.

\paragraph{Frozen prospective replacement.}
To prevent this defect from recurring, the artifact includes
\artifact{scripts/make\_prospective\_validation\_sample.py} and
\artifact{corpus/annotations/prospective\_validation\_manifest.json}. Before
any outcome was inspected, seed \slug{20260728} selected 24 repositories
uniformly per framework and one eligible corrected-v2 file uniformly within
each selected repository: 96 files total, with exact inclusion probabilities
and pinned source hashes. A separately seeded 8-per-framework subset (32 total)
is locked for dual-human graph reconstruction. The complete staging,
adjudication, and stop/go rules are in
\artifact{PROSPECTIVE\_VALIDATION\_PROTOCOL.md}. These prospective records
contain no labels and contribute no number to the present results; they replace
the historical quantities only after two independent annotators complete the
protocol.

\begin{table}[htbp]
\centering\small
\caption{Achieved sampling fractions (descriptive, not design). Source:
\artifact{corpus/real\_world/metadata\_v2.json} and
\artifact{revision\_analyses.json}.}
\label{tab:fractions}
\begin{tabular}{@{}lrrr@{}}
\toprule
\textbf{Framework} & \textbf{Corpus $N$} & \textbf{Sample $n$} & \textbf{$n/N$} \\
\midrule
LangGraph & 400 & 40 & 10.0\% \\
AutoGen   & 359 & 35 & 9.7\% \\
CrewAI    & 113 & 20 & 17.7\% \\
ADK       & 50  & 24 & \textbf{48.0\%} \\
\midrule
\textbf{Total} & \textbf{922} & \textbf{119} & 12.9\% \\
\bottomrule
\end{tabular}
\end{table}

\subsection{Why inclusion probabilities are undefined}

Four independent reasons, each sufficient alone.

\begin{enumerate}
\item \textbf{The first-stage selection is a proprietary ranking, not a
randomization.} \slug{gh search code} with no \slug{--sort} returns GitHub's
``best match'' order: undisclosed, non-stationary, and deterministic given the
index. Truncating it at 40 gives every candidate an inclusion \emph{indicator} of
0 or 1, not a probability. There is no $\pi_i$ to compute and no way to bound it.
\item \textbf{The frame's denominator is unknown.} GitHub does not publish the
size or membership of the code-search index, and no total-hit count was recorded
at mining time---\slug{gh\_search} (\artifact{mine\_github\_gh.py}:52--70)
discards everything except the capped item list. Even a Horvitz--Thompson
estimator with guessed weights has no $N$ to expand to.
\item \textbf{The second stage compounds the first deterministically.}
\slug{unique[:180]} is a hard cut on the same rank order---an entire framework
(AutoGen, ADK) fell below it in run~1. That is a selection mechanism correlated
with the ranking signal, not with anything exchangeable.
\item \textbf{The third stage is a census with a non-random thinning.} Within a
surviving repo all \slug{*.py} files are taken (cluster sizes 1 to 55), and the
slug-collision overwrite then deletes 150 of them by filename coincidence.
\end{enumerate}

Add the unreconstructible hand-selection of the 119 and there is no level of the
design at which a probability model can be attached.

\paragraph{Consequence.} No design-based variance estimator is licensed. The
Wilson intervals and the repository-clustered bootstrap are legitimate as
\emph{model-based uncertainty for the 119 observations under an
i.i.d.-within-stratum assumption}, and are the right thing to report, but they
quantify \textbf{only} sampling noise \textbf{within the corpus} and carry
\textbf{no} coverage guarantee for GitHub. The dominant error term---selection
bias of the corpus itself---is unquantified and unquantifiable from these
artifacts. For the same reason the finite-population correction an earlier draft
applied (\artifact{scripts/reviewer\_analyses.py}:148, \slug{fpc = 1 - n/N\_f})
is \emph{not licensed}, since an FPC presupposes simple random sampling without
replacement from $N_f$; the clustered-bootstrap interval is the defensible one
and is the one quoted.

\subsection{The reweighting: what it is and is not}

\artifact{scripts/reviewer\_analyses.py}:326--372 post-stratifies each
per-framework rate by corpus framework shares,
$\hat p = \sum_f (N_f/N)\,(d_f/n_f)$, with CIs from a repository-clustered
bootstrap within strata at fixed weights. The correct name for this quantity is a
\textbf{corpus-share-weighted descriptive statistic}: the rate the 922-graph
corpus would show \emph{if} each framework's 119-sample rate held across that
framework's corpus members. It is a within-corpus standardization that removes
the ADK over-representation of Table~\ref{tab:fractions}. It is \textbf{not} a
prevalence estimate for GitHub, and it is \textbf{not} post-stratification in the
survey sense, which requires known \emph{population} control totals; here the
control totals are the corpus's own counts, so the procedure cannot correct any
bias introduced before the corpus existed---which is all of the bias.

\subsection{The corpus framework mix, and a retracted ``correction''}
\label{sec:census}

The authoritative framework mix is obtained by counting the \slug{framework}
field of each of the 922 non-self-repo graph JSONs under
\artifact{corpus/real\_world/graphs/} directly---the field the extractor writes
from the file it actually parsed. That count is:

\begin{table}[htbp]
\centering\small
\caption{Corpus framework mix. Source: direct count of the \slug{framework}
field over all 922 non-self-repo files in \artifact{corpus/real\_world/graphs/};
independently confirmed by \artifact{matched\_fidelity.json} $\to$
\slug{drops.corpus\_framework\_distribution} and by the hard-coded weights at
\artifact{scripts/reviewer\_analyses.py}:57.}
\label{tab:census}
\begin{tabular}{@{}lrr@{}}
\toprule
\textbf{Framework} & \textbf{$N$} & \textbf{Share} \\
\midrule
LangGraph & 400 & 43.4\% \\
AutoGen   & 359 & 38.9\% \\
CrewAI    & 113 & 12.3\% \\
ADK       & 50  & 5.4\% \\
\midrule
\textbf{Total} & \textbf{922} & 100\% \\
\bottomrule
\end{tabular}
\end{table}

\paragraph{A retraction.}
An earlier revision of our own methods audit
(\artifact{papers/paper1/aaai/SAMPLING\_DESIGN.md}, \S6.2) asserted that this mix
was wrong and that the artifacts implied AutoGen~357 / CrewAI~115. \textbf{That
assertion was itself the error, and we retract it here.} It was produced by
reading framework labels for the \emph{v1} corpus out of the \emph{v2} re-mine
records---the defect at \artifact{scripts/matched\_fidelity.py}:251 documented in
Section~\ref{sec:fwbug}---which mislabels exactly the seven slug-collision cases
in which the two miners resolved the same slug to different source files. The
net effect of those seven, $\text{AutoGen}\,{-}2$ and $\text{CrewAI}\,{+}2$, is
precisely the discrepancy the audit reported and misattributed to a self-repo
bookkeeping slip.

\paragraph{No published number changes.}
We checked the downstream consequence rather than assuming it. The
post-stratification weights hard-coded at
\artifact{scripts/reviewer\_analyses.py}:57 are
$\{\text{langgraph}\,400,\ \text{crewai}\,113,\ \text{autogen}\,359,\
\text{adk}\,50\}$---the \emph{correct} mix---and
\artifact{corpus/real\_world/reviewer\_analyses.json} $\to$
\slug{2\_post\_stratification.corpus\_composition} records the same values as
actually used. Every post-stratified figure was therefore computed with the
right weights throughout, and none of them moves---neither the superseded
two-check row (structural $0.23\%$, policy $1.92\%$, composite $2.15\%$) nor the
HGP-1 figures the main text's Table~2 now reports (structural $0.23\%$, policy
$12.30\%$, composite $12.52\%$), which use the identical weights. The erroneous $357/115$ pair
never propagated beyond the prose of the audit document and the corresponding
row of an earlier draft of this supplement, both now corrected. The sampling
fractions in Table~\ref{tab:fractions} are the only numbers in this supplement
that changed, and they change in the fourth significant figure
(AutoGen $9.8\%\to9.7\%$, CrewAI $17.4\%\to17.7\%$).

\subsection{What the design does support}

Descriptive claims about this corpus, with full provenance (every graph is
pinned to \slug{(repo, SHA, path)} and re-extractable via
\artifact{scripts/remine\_pinned.py}); instrument-fidelity claims, which are
properties of the instrument on these inputs and do not require a probability
sample; the existence and direction findings (a biased frame makes a \emph{low}
observed defect rate more notable, not less, provided no population number is
attached); and relative comparisons within the corpus. What it does \emph{not}
support is any sentence of the form ``$X\%$ of agent workflows on GitHub''.
Finally, the frame itself is \textbf{not reproducible}---re-running the 11
queries today returns a different top-40 each. What \emph{is} reproducible is the
corpus \emph{given} \artifact{candidates.json}/\artifact{repos*.json}.
Reproducibility claims should be scoped to ``the corpus is reproducible from the
pinned manifest'', never to ``the mining is reproducible''.

% ===========================================================================
\section{The Prospectively Specified Human-Gate Policy (HGP-1)}
\label{sec:hgp}

Included because it is the instrument behind the human-gate audit. The following inlines
\artifact{papers/paper1/aaai/HUMAN\_GATE\_POLICY.md}. Machine-readable record:
\artifact{corpus/real\_world/human\_gate\_audit.json}; script:
\artifact{scripts/human\_gate\_audit.py}.

\paragraph{Status and purpose.}
Prospectively specified: written before any of the 32 candidate workflows were
assessed for gate presence, and before any comparison against
\slug{human\_presence} or \slug{human\_gate\_coverage} output. Its purpose is
to replace a conflated estimand---three cases produced by two different checks
over two different artifact populations---with one operational procedure applied
uniformly to every workflow in the sample that contains a side-effecting tool at
\emph{source} level, so that false negatives outside both flag universes become
detectable.

\paragraph{Unit and evidence base.}
One workflow $=$ one mined source file (a slug in
\artifact{corpus/real\_world/gt\_triage\_results.json}); denominator $n=119$.
Evidence is the source text under \artifact{corpus/real\_world/sources/}, or the
recorded \slug{repo}/\slug{sha}/\slug{path}. Extracted graphs may be
consulted as a \emph{navigation aid only}: a verdict must be justified by a
source line, never by a graph node kind, because the
\slug{tool}$\rightarrow$\slug{llm} error is precisely the failure under audit.

\begin{table}[htbp]
\centering\small
\caption{(a) What counts as a side-effecting / irreversible tool. Classification
is by what the callee actually does in source, not by its name.}
\label{tab:hgp_a}
\begin{tabular}{@{}lp{0.66\linewidth}@{}}
\toprule
\textbf{Category} & \textbf{Definition and positive examples} \\
\midrule
\slug{destructive} & Deletes, drops, truncates, overwrites, or irreversibly
removes existing persistent state: \slug{DROP TABLE}, \slug{DELETE FROM},
\slug{os.remove}, \slug{shutil.rmtree}, \slug{git reset --hard}, cache
purge, resource teardown. \\
\slug{external\_write} & Creates or mutates persistent state outside the
process---databases, cloud storage, repositories, ticketing/CRM---or a
non-idempotent HTTP \slug{POST}/\slug{PUT}/\slug{PATCH}/\slug{DELETE} to
a third-party API: \slug{INSERT}/\slug{UPDATE}/\slug{MERGE},
\slug{open(path,'w')} outside a temp dir, \slug{bigquery.insert\_rows},
\slug{firestore.set}, \slug{git commit}/\slug{push}, issue/PR creation. \\
\slug{exec} & Executes attacker- or model-influenced code or commands:
\slug{subprocess.run}, \slug{os.system}, \slug{exec}/\slug{eval},
\slug{PythonAstREPLTool}, \slug{PythonCodeExecutionTool},
\slug{kubectl apply}, container exec, \slug{docker run} on a shared daemon. \\
\slug{financial} & Moves money or creates a financial obligation: payment
capture, transfer, trade order, invoice issue, refund, charge. \\
\slug{comms\_third\_party} & Delivers a message outside the operator's trust
boundary: SMTP send, SendGrid/Gmail send, SMS, Slack/Discord channel post,
social-media publish, outbound customer webhook. \\
\bottomrule
\end{tabular}
\end{table}

\paragraph{(b) What counts as a human gate.}
A control point at which execution \emph{blocks pending a human decision that can
prevent the side effect}. Exactly one of: (1) an explicit human checkpoint
node---graph node of kind \textsc{human}, LangGraph \slug{interrupt()},
\slug{interrupt\_before=[n]} or \slug{NodeInterrupt}, ADK
\slug{LongRunningFunctionTool} or tool-confirmation callback, CrewAI task with
\slug{human\_input=True}; (2) an equivalent confirmation step in the node
body---a blocking read of human input (\slug{input()}, \slug{click.confirm()},
a Streamlit/Gradio button handler, an awaited approval message) \emph{whose
negative outcome causes the side-effecting call to be skipped}, with the negative
branch present in source; or (3) AutoGen \slug{UserProxyAgent} with
\slug{human\_input\_mode="ALWAYS"} positioned so the human speaks between the
proposing agent and the executing agent.

\emph{Explicitly not a gate} (fixed before case labeling): a
generic conversational REPL in which the human supplies the \emph{task} and the
agent then autonomously chooses and executes tools (the human authorises a goal,
not the specific irreversible action); \slug{human\_input\_mode="NEVER"} or
\slug{"TERMINATE"}; logging, printing, or streaming with no blocking decision;
post-hoc review occurring \emph{after} the side effect has landed; and a gate on
some other branch that does not lie on the path to the side-effecting call.

\paragraph{(c) Violation condition.}
A workflow is a violation iff there exists at least one path from the workflow
entry point to an invocation of a side-effecting tool that contains no human
gate. Equivalently: the set of human gates does not \emph{dominate} the
side-effecting call site. A gate on one branch does not discharge an ungated
sibling branch, and one ungated reachable side-effecting tool suffices regardless
of how many others are correctly gated.

\paragraph{(d) Exclusions.}
E1 read-only tools (retrieval, search, \slug{SELECT}, HTTP \slug{GET}, vector
query, file read, web scrape); E2 pure computation; E3 in-process state only
(Python dict, agent/session state, LangGraph channel state, conversation
history); E4 LLM text generation, unless the text is then persisted or
transmitted by a separate call; E5 genuine mocks and stubs, the exclusion
attaching to the \emph{implementation} not the file's name; E6 constructively
sandboxed side effects (writes confined to \slug{tempfile}/\slug{tmp\_path},
or execution inside a container created and torn down for that call), with the
confinement visible in source.

\emph{Explicitly not excluded} (fixed before case labeling): N1 demo/toy/tutorial/example
status---``it's only a demo'' is a statement about intent, not behaviour, and
intent is recorded separately without changing the verdict; N2 test files, which
require E5 or E6 to be exempt; N3 unconfigured credentials---a send that would
fail for want of an API key still counts, since the missing gate is the defect;
N4 local filesystem writes outside a temp directory.

\paragraph{(e) Undecidable cases.}
The verdict domain is exactly
$\{$\slug{violation}, \slug{compliant}, \slug{arguable},
\slug{source\_unavailable}$\}$. \slug{arguable} is mandatory when (i) the
side-effecting callee is imported from a module absent from the corpus and its
behaviour cannot be established, (ii) tool dispatch is fully dynamic, or
(iii) a candidate gate exists but whether it dominates the call site depends on
runtime configuration not fixed in source; an \slug{arguable} verdict must
state which applies. \textbf{An undecidable case is never recorded as
\slug{compliant}}---silent downgrading of uncertainty to a negative is the
specific bias this policy exists to prevent. We report
$\text{violation}/119$ with a descriptive Wilson interval, plus a sensitivity band whose lower
arm counts only \slug{violation} and whose upper arm counts
$\text{\slug{violation}}+\text{\slug{arguable}}$; both are reported regardless of direction.

\subsection{Applied results}

All 32 sources were read and none dropped;
\slug{Sidreyas\_\_The\_Grand\_AI\_Repo\_\_builder} was missing from
\artifact{corpus/real\_world/sources/} and was recovered from GitHub at the
recorded sha \slug{61bb48c6}, so there are \textbf{zero}
\slug{source\_unavailable} cases. One threshold rule was clarified at
application time and applied uniformly in both directions: a side-effecting
callable counts only if it is agent/LLM-invocable or is called inside a graph
node body that executes during a workflow run; script-level setup, teardown, and
reporting code is outside the threshold.

\begin{table}[htbp]
\centering\small
\caption{HGP-1 outcome over the 32 source-level side-effecting workflows.
Source: \artifact{corpus/real\_world/human\_gate\_audit.json}.}
\label{tab:hgp_out}
\begin{tabular}{@{}lr@{}}
\toprule
\textbf{Verdict} & \textbf{Count} \\
\midrule
\slug{violation}          & \textbf{12} \\
\slug{compliant}          & 9 \\
\slug{arguable}           & 11 \\
\slug{source\_unavailable} & 0 \\
\bottomrule
\end{tabular}
\end{table}

\begin{itemize}
\item Prevalence over the sample: $\mathbf{12/119 = 10.08\%}$, Wilson 95\%
$[5.86\%, 16.80\%]$.
\item Sensitivity upper arm ($\text{\slug{violation}}+\text{\slug{arguable}} = 23/119$):
$[13.24\%, 27.34\%]$.
\item Prevalence among the audited side-effecting workflows: $12/32 = 37.50\%$
$[22.93\%, 54.75\%]$.
\end{itemize}

\paragraph{HGP-1 supersedes the earlier two-check procedure.}
HGP-1 is the paper's \textbf{primary} human-gate instrument, and the main text
and its Table~2 report the $12/119$ ($10.1\%$) $[5.86, 16.80]$ figure derived
here. It \emph{supersedes} the $3/119$ ($2.52\%$) figure an earlier draft
reported: of those 3 original cases, HGP-1 retains 1 and finds 11 additional
violations, a fourfold upward revision that places the superseded figure
\emph{outside} the new interval.

The supersession is a matter of instrument quality, not of preferring the larger
number. The earlier figure was produced by two \emph{different} checks
(\slug{human\_presence} and \slug{human\_gate\_coverage}) applied to two
\emph{different} artifact populations (mined graphs and reconstructed graphs),
so it was never a single estimand; and both checks read gate presence off a
graph node's \emph{kind}, which is exactly the field the
\slug{tool}$\rightarrow$\slug{llm} error corrupts. HGP-1 replaces this with one
prospectively specified procedure applied uniformly to source text, with a verdict
domain that forbids recording an undecidable case as compliant. The two
subsections below---the three violations invisible to \emph{both} earlier
checks, and the eight non-discharging ``gates'' that both checks accepted---are
the direct evidence that the superseded procedure was measuring gate
\emph{labels} rather than gate \emph{behaviour}.

\paragraph{False negatives.}
Three genuine violations lie outside \emph{both} flag universes---neither
\slug{human\_presence} nor \slug{human\_gate\_coverage} fires on them:
\slug{Arhamirfan\_\_Snake-Agents\_\_snake\_autogen},
\slug{Itzkuldeep\_\_AgentE-com\_\_main}, and
\slug{SAMAR-CODE404\_\_backend\_\_Main}. A further 9 violations are missed by
the risk-aware check specifically (11 of the 12 have
\slug{human\_gate\_coverage\_failed = false}); those 9 fall inside the
\slug{human\_presence} universe, but that universe contains 106 of 119
workflows and so carries almost no information.

\subsection{Placebo gates: a distinct failure mode}

Of the 32, \textbf{8} carry a human gate that cannot discharge the obligation,
and 5 are outright placebos---a node typed \slug{human} in the \emph{reference}
graph whose body can never block execution.

\begin{table}[htbp]
\centering\small
\caption{The 8 non-discharging gates. Source:
\artifact{corpus/real\_world/human\_gate\_audit.json}.}
\label{tab:placebo}
\begin{tabular}{@{}p{0.36\linewidth}p{0.56\linewidth}@{}}
\toprule
\textbf{Slug} & \textbf{Nature of the gate} \\
\midrule
\slug{SAMAR-CODE404\_\_backend\_\_Main} & placebo: 6 \slug{*\_human\_approval}
nodes whose body is \slug{proceed = self.approval} against a constructor flag
hard-set \slug{True} \\
\slug{Arhamirfan\_\_Snake-Agents\_\_snake\_autogen} & placebo:
\slug{UserProxyAgent} with \slug{human\_input\_mode="TERMINATE"} \\
\slug{timleow\_\_gym-kfc-daddies\_\_cli\_ui} & placebo:
\slug{get\_user\_prompt} is a task REPL, not an approval \\
\slug{Sidreyas\_\_The\_Grand\_AI\_Repo\_\_builder} & inert:
\slug{web\_user\_proxy} serialized into a config, never executed \\
\slug{waqasniazi9\_\_voice-agent\_\_builder} & inert: \slug{web\_user\_proxy}
serialized into a config, never executed \\
\slug{Itzkuldeep\_\_AgentE-com\_\_main} & non-dominating: real \slug{ALWAYS}
proxy, but ordered after the executing agent \\
\slug{Sujas-Aggarwal\_\_langraph-chatbot\_\_v2.6.0} & partial:
\slug{confirmation\_node} auto-confirms above a similarity threshold \\
\slug{merdandt\_\_SalesShortcut\_\_websiter\_creator\_agent} & unverifiable:
\slug{request\_URL} typed \slug{human}, body absent \\
\bottomrule
\end{tabular}
\end{table}

\textbf{All 7} side-effecting workflows that both existing checks cleared were
cleared because the reference graph contained a node typed \slug{human}; on
source inspection \textbf{none of the 7 is a genuine dominating gate}. The bias
is therefore \emph{not} confined to the extractor: the source-reconstructed
reference graphs themselves overstate human oversight, because \slug{human} is assigned
from node name or framework type rather than from whether the body can block. Any
estimand computed over those graphs---including the corrected, risk-aware
one---is biased \textbf{downward}, systematically rather than randomly. The
\slug{human\_detection} fidelity metric compounds this: for
\slug{Arhamirfan\_\_Snake-Agents\_\_snake\_autogen} and
\slug{Itzkuldeep\_\_AgentE-com\_\_main} it records $\text{tp}=1$,
$\text{recall}=1.0$, scoring the extractor as \emph{perfectly} recovering human
oversight that does not exist; for \slug{SAMAR-CODE404\_\_backend\_\_Main} it
records $\text{fn}=6$, penalising the extractor for failing to reproduce six
placebo gates.

\paragraph{Honest limitations of HGP-1.}
The 11 \slug{arguable} cases are dominated by reason (i): the side-effecting
callee lives in a module the mining pipeline never captured (single-file
snapshots of multi-file packages). A repository-level corpus would resolve most
of them, and several look likely to resolve toward \slug{violation}---most
notably \slug{Zen7-Labs\_\_Zen7-Payment-Agent\_\_agent}, a crypto
payment/settlement pipeline whose prompts forbid confirmation four separate times
and where only the callee body, not the missing gate, is unverifiable. They are
held at \slug{arguable} because the policy requires it, not because the evidence
is balanced. The 87 workflows screened as having no source-level side effect
were not re-read under HGP-1, so upstream screening error can change the
$12/119$ count in either direction. Finally the audit is
\textbf{single-rater}; \slug{ed-donner\_\_agents\_\_sequential\_agents}
(prompt-level confirmation only) and \slug{microsoft\_\_ACV\_\_test\_cache\_agent}
(compliant) are the two calls most likely to move under a second rater.

% ===========================================================================
\section{The Slug-Collision Keying Defect and the Matched-Set Methodology}
\label{sec:collisions}

\emph{Elaborates} \artifact{sections/03\_realworld.tex}, paragraph
\emph{A keying defect, and how we corrected for it}. All numbers here come from
\artifact{corpus/real\_world/matched\_fidelity.json} $\to$
\slug{provenance\_collisions}.

\subsection{The defect}

Graphs are keyed by
$\text{\slug{slug}} = \langle\text{repo}\rangle\_\_\langle\text{file basename}\rangle$,
which is not unique within a repository. The artifact's own note states the
consequence precisely: \slug{remine\_pinned.load\_records()} resolves
collisions \emph{first-wins}, so the v2 instrument may re-mine a \emph{different
file} than v1 did for the same slug. Such a graph pair compares two different
programs, not two instrument versions. \textbf{76 slugs collide} in total; 51 of
them are in the matched corpus and 9 are in the matched fidelity set.

\begin{table}[htbp]
\centering\small
\caption{The nine colliding slugs inside the matched fidelity set, with the
number of distinct source files that mapped to each. Source:
\artifact{matched\_fidelity.json} $\to$
\slug{provenance\_collisions.fidelity\_collision\_slugs}.}
\label{tab:collisionslugs}
\begin{tabular}{@{}lr@{}}
\toprule
\textbf{Slug} & \textbf{Files collapsed} \\
\midrule
\slug{svpino\_\_ml.school\_\_agent}                        & 5 \\
\slug{Saimoguloju\_\_AI-Agents\_\_sample\_agent}            & 4 \\
\slug{lordpython\_\_multi-agent-video-system\_\_agent}      & 4 \\
\slug{Astrodevil\_\_ADK-Agent-Examples\_\_agent}            & 3 \\
\slug{MikhailMostWanted\_\_Astra\_\_app\_team}              & 3 \\
\slug{joansantoso\_\_iox-adk\_\_agent}                      & 2 \\
\slug{shazy89\_\_ai-learning\_\_agent}                      & 2 \\
\slug{volcengine\_\_veadk-python\_\_main}                   & 2 \\
\slug{winstonbartlegod\_\_enhanced-multimodal-rag\_\_v1}    & 2 \\
\bottomrule
\end{tabular}
\end{table}

Table~\ref{tab:collisionslugs} lists them. The pattern is visible in the names:
repositories that keep one \slug{agent.py} per framework subdirectory.

\subsection{Why zero of the 58 pairs are recoverable}

\begin{table}[htbp]
\centering\small
\caption{Re-keying recovery analysis. Source:
\artifact{matched\_fidelity.json} $\to$
\slug{provenance\_collisions.rekey\_recovery} (the artifact also enumerates all
58 unrecoverable records individually under \slug{unrecoverable\_detail}).}
\label{tab:rekey}
\begin{tabular}{@{}lr@{}}
\toprule
\textbf{Quantity} & \textbf{Count} \\
\midrule
Colliding slugs                                   & 76 \\
Present on disk, v1                               & 58 \\
Present on disk, v2                               & 51 \\
v2 records re-keyable to an exact file path       & 51 \\
v1 records re-keyable, even only by framework     & 5 \\
\midrule
\textbf{Usable v1--v2 pairs recovered}             & \textbf{0} \\
\textbf{Must drop}                                 & \textbf{58} \\
\bottomrule
\end{tabular}
\end{table}

The asymmetry in Table~\ref{tab:rekey} is the whole explanation. The v2 miner
records enough provenance that all 51 of its colliding records can be re-keyed to
an exact file path. The v1 miner resolved collisions \emph{last-wins without
recording which file survived}, so v1's file identity is \emph{structurally}
unrecoverable---at best 5 of 58 can be pinned even to a framework. A pair
requires both halves; since the v1 half is missing for at least 53 of 58 and
under-determined for the rest, \textbf{zero usable pairs} can be reconstructed.
This is not a matter of effort: the information was never written down. All 58
are dropped, and the released instrument fixes the keying.

\subsection{The drop is not uniform: ADK-concentrated bias}

\begin{table}[htbp]
\centering\small
\caption{Framework composition of the fidelity sets before and after the
collision exclusion. Source: \artifact{matched\_fidelity.json} $\to$
\slug{provenance\_collisions.fidelity\_drop\_bias}.}
\label{tab:dropbias}
\begin{tabular}{@{}lrrrr@{}}
\toprule
\textbf{Framework} & \textbf{Matched (115)} & \textbf{Dropped (9)} &
\textbf{Remaining (106)} & \textbf{Drop rate} \\
\midrule
ADK       & 23 & 5 & 18 & \textbf{21.7\%} \\
AutoGen   & 32 & 2 & 30 & 6.2\% \\
LangGraph & 40 & 2 & 38 & 5.0\% \\
CrewAI    & 20 & 0 & 20 & 0.0\% \\
\midrule
\textbf{Total} & \textbf{115} & \textbf{9} & \textbf{106} & 7.8\% \\
\bottomrule
\end{tabular}
\end{table}

ADK loses more than a fifth of its stratum while CrewAI loses none. The cause is
mechanical rather than statistical: ADK repositories disproportionately use
per-framework directories containing identically named \slug{agent.py} files,
which is exactly the collision pattern. \textbf{ADK's row in the main paper's
fidelity table is therefore the least stable and should be read as such}---a
judgment that Table~\ref{tab:fidelity_v2}'s interval, $[0.067, 0.636]$ on edge
recall, independently confirms.

\begin{table}[htbp]
\centering\small
\caption{Reference-graph confidence of the dropped versus retained graphs.
Source: \artifact{matched\_fidelity.json} $\to$
\slug{provenance\_collisions.fidelity\_drop\_bias}.}
\label{tab:dropconf}
\begin{tabular}{@{}lrrr@{}}
\toprule
& \textbf{high} & \textbf{medium} & \textbf{low} \\
\midrule
Dropped (9)  & 8 (88.9\%) & 1 (11.1\%) & 0 \\
Matched (115) & 75 (65.2\%) & 32 (27.8\%) & 8 (7.0\%) \\
\bottomrule
\end{tabular}
\end{table}

The exclusion also removes proportionally \emph{more} high-confidence reference
graphs than low ($88.9\%$ of the dropped set versus $65.2\%$ of the matched set),
so the retained set is slightly lower-quality on average than the set it came
from. We report the direction because it is adverse to us; we do not attempt to
correct for it, since with $n=9$ any correction would be noise.

\subsection{Effect on the structural-flag comparison}

\begin{table}[htbp]
\centering\small
\caption{Structural flags on the matched corpus, with and without the colliding
slugs. Source: \artifact{matched\_fidelity.json} $\to$
\slug{matched\_structural\_flags} and its \slug{collision\_free} sub-key.}
\label{tab:structflags}
\begin{tabular}{@{}lrr@{}}
\toprule
\textbf{Set} & \textbf{v1 flagged} & \textbf{v2 flagged} \\
\midrule
Matched corpus ($n=904$)              & 313 & 239 \\
Collision-free ($n=853$, 51 removed)  & 299 & 229 \\
\bottomrule
\end{tabular}
\end{table}

The main text quotes the collision-free row ($299\to229$), which is the
defensible one. For completeness: 75 workflows are flagged \emph{only} in v1 and
1 \emph{only} in v2, and of the 150 graphs whose structure changed between
instruments, 40 changed because of a collision (different file) and 110 because
of a genuine extractor change---the artifact separates these under
\slug{corpus\_structurally\_changed}, which is what licenses the main text's
causal reading of the corrected-instrument ablation.

\subsection{A second defect with the same root cause: framework mislabelling}
\label{sec:fwbug}

The collision defect has a downstream sibling that we found while preparing this
supplement, and it is worth stating separately because it bit our own audit
rather than the paper.

\artifact{scripts/matched\_fidelity.py} built the corpus framework census with
\begin{quote}
\slug{corpus\_fw = Counter(status.get(s, \{\}).get("framework") for s in corpus\_v1)}
\end{quote}
at line~251, where \slug{status} is the map of \emph{v2} re-mine records
(\slug{load\_v2\_status}) but \slug{corpus\_v1} is the set of \emph{v1} slugs. It
therefore labels every v1 graph with the framework of whatever file \emph{v2}
resolved that slug to. For a non-colliding slug the two agree and the bug is
invisible; for a colliding slug they need not, because v1 resolved last-wins and
v2 resolves first-wins---the very asymmetry of
Section~\ref{sec:collisions}. \textbf{Seven slugs disagree, and all seven are
slug collisions:}

\begin{table}[htbp]
\centering\small
\caption{The seven slugs whose framework label differs between the v1 graph
JSON's own \slug{framework} field and the v2 re-mine record. All seven appear in
the 76-slug colliding population. Source: recount over
\artifact{corpus/real\_world/graphs/} versus
\artifact{corpus/real\_world/metadata\_v2.json}.}
\label{tab:fwbug}
\begin{tabular}{@{}p{0.46\linewidth}ll@{}}
\toprule
\textbf{Slug} & \textbf{v1 graph} & \textbf{v2 record} \\
\midrule
\slug{Narwhal-Lab\_\_MagicSkills\_\_model}          & crewai    & autogen \\
\slug{Saimoguloju\_\_AI-Agents\_\_sample\_agent}     & autogen   & langgraph \\
\slug{ed-donner\_\_agents\_\_agent}                  & langgraph & crewai \\
\slug{ed-donner\_\_agents\_\_main}                   & langgraph & crewai \\
\slug{feixiao\_\_ai\_\_chap06}                       & autogen   & crewai \\
\slug{feixiao\_\_ai\_\_chap07}                       & autogen   & crewai \\
\slug{martimfasantos\_\_ai-agents-frameworks\_\_00\_hello\_world} & crewai & langgraph \\
\bottomrule
\end{tabular}
\end{table}

The net displacement is $\text{AutoGen}\,{-}2$, $\text{CrewAI}\,{+}2$, which is
exactly the spurious $359/113 \to 357/115$ shift that Section~\ref{sec:census}
retracts. \textbf{Blast radius, checked rather than assumed:} the bug touches
only this one \slug{Counter}. The fidelity sets take their framework labels from
a different map (built from the validation sample's own records), so every
number in Sections~\ref{sec:fidelity} and~\ref{sec:collisions}---the $n=106$,
$n=115$ and $n=119$ fidelity tables, the drop-bias table, the re-keying
analysis, and the structural-flag counts---is \emph{bit-identical} before and
after the fix, which we verified by re-reading the regenerated artifact
field-by-field. Five of the seven mislabelled slugs are not in the 119-graph
validation sample at all, and the other two are among the 9 colliding graphs
already excluded from the $n=106$ set. The fix is committed and the artifact
regenerated.

\begin{remark}[Why we report this]
Both defects are instances of one failure: a non-unique key
($\langle\text{repo}\rangle\_\_\langle\text{basename}\rangle$) joined across two
datasets that resolve it by opposite tie-break rules. The first cost us 58
comparison pairs; the second cost us a false ``correction'' in our own audit that
we then propagated into a draft of this supplement. Neither was caught by a
test, because both produce plausible numbers that sum correctly.
\end{remark}

\subsection{Corpus-level drops, for completeness}

Independently of collisions, 18 of the 922 v1 graphs have no v2 counterpart:
13 because the repository became unreachable at re-mining time and 5 because
re-extraction produced no graph. By framework: AutoGen 14, ADK 2, LangGraph 1,
CrewAI 1; 13 of the 18 come from the single repository
\slug{Sidreyas/The\_Grand\_AI\_Repo}. This yields the matched corpus of
$n=904$. There are \textbf{zero} orphan v2 graphs absent from v1. Source:
\artifact{matched\_fidelity.json} $\to$ \slug{drops}.

% ===========================================================================
\section{Positioning Against Related Measurement and Enforcement Work}
\label{sec:related}

\emph{Referenced from} \artifact{sections/05\_related.tex}, which states that the
supplement places these efforts along six axes. Table~\ref{tab:related} is that
placement; it was moved here from the main paper to meet the page limit.

The six axes are \textbf{corpus} (what was studied, and how much of it),
\textbf{frameworks} (breadth of agent frameworks covered), \textbf{model} (the
formal object, if any, that analysis runs over), \textbf{properties} (what is
checked or catalogued), \textbf{validation} (how the findings were checked), and
\textbf{guarantees} (what is claimed to follow). The split the table makes
visible is that prior measurement studies characterise \emph{which} defects
occur, and prior enforcement systems address \emph{how} to constrain agents at
runtime, while neither line measures the fidelity of the extracted model on
which any static check depends.

\begin{table}[htbp]
\centering\footnotesize
\setlength{\tabcolsep}{4pt}
\caption{Positioning against measurement and enforcement work for LLM agents,
along the six axes of Section~\ref{sec:related}. Ours measures the fidelity of
the extracted model on which static checks depend, separates distinct audit
quantities, and makes soundness conditional on caller-asserted event coverage.}
\label{tab:related}
\resizebox{\textwidth}{!}{%
\begin{tabular}{@{}lllllll@{}}
\toprule
Work & Corpus & Frameworks & Model & Properties & Validation & Guarantees \\
\midrule
\citet{ning2024agentable}    & dev discussions      & agnostic        & ---               & defect taxonomy       & manual        & none \\
\citet{xue2025bugs}          & ${\sim}1{,}026$ bugs & orchestration   & ---               & 9 root causes         & manual        & none \\
\citet{shen2025securitydebt} & $221$ vulns          & agent apps      & ---               & 14 vuln.\ types       & manual        & none \\
\citet{wang2026agentflow}    & $5{,}399$ progs      & cross-framework & dep.\ graph       & security patterns     & ---           & none (unvalidated) \\
\citet{wang2025agentspec}    & case studies         & runtime         & rule DSL          & safety rules          & case studies  & runtime enforcement \\
\citet{kamath2026agentc}     & case studies         & generation-time & temporal aut.\    & temporal constraints  & conformance   & conformance (runtime) \\
\midrule
\textbf{Ours} & $922$ graphs, $252$ repos & $4$ (cross) & graph\,+\,LTL$_f$ &
struct.\,+\,policy quantities & LLM-recon.; \humanValidationStatus &
caller-asserted soundness (rel.\ authored abstraction) \\
\bottomrule
\end{tabular}}
\end{table}

\paragraph{On our own ``Guarantees'' entry.}
The wording is deliberately weaker than an unqualified soundness claim, and is
the same claim Section~\ref{sec:proof} proves. What the system delivers is
soundness \emph{relative to an authored abstraction} whose event coverage the
caller asserts; exact provenance qualifies witnesses but never discharges that
assertion (Proposition~\ref{prop:gap}). Reading the cell
as ``verified soundness'' would overstate the result in exactly the way
Section~\ref{sec:h1} warns against. Two neighbouring cells are similarly
hedged on purpose: our ``Validation'' entry records that the reference graphs
are LLM reconstructions and reports the status of independent human validation,
and our ``Corpus'' entry counts \emph{extracted file-level graphs}, not
executable workflow roots or deployed systems---so it is not directly
comparable to a bug count or a vulnerability count in the rows above.

\paragraph{What the table does not encode.}
It is a capability placement, not a benchmark: no row was re-run by us, and the
corpus sizes are not commensurable across rows (bug reports, vulnerability
records, programs, and extracted graphs are different units). We make no claim
of the form ``ours is $N\times$ larger''.

\section{Independent Human Validation}
\label{sec:human-validation}

\ifnum\humanNflags>0
Two independent annotators labeled \humanNflags{} defect items. Human--human
raw agreement is \humanRaw{}, Cohen's $\kappa=\humanKappa$, and
Krippendorff's $\alpha=\humanAlpha$. Cohen's $\kappa$ against the committed LLM
labels is \humanLLMkappaA{} for annotator A and \humanLLMkappaB{} for annotator
B. The annotators also independently reconstructed \humanGraphN{} locked
workflows: human--human mean edge F1 is \humanGraphEdgeF{} and node-kind
agreement is \humanGraphKind{}; annotator A's mean edge F1 against the
LLM-reconstructed references is \humanRefEdgeF{}. Full confusion matrices,
repository-cluster bootstrap intervals, per-item graph scores, and completion
counts are in \artifact{corpus/annotations/human\_agreement.json}.
\else
Two-annotator validation has not yet been executed. The locked samples,
independent worksheets, annotation guide, and scoring script are included under
\artifact{corpus/annotations/} and
\artifact{scripts/human\_agreement.py}. Running the latter with
\slug{--latex} updates this section and the main paper from the same generated
macro file. Until then, LLM-assisted labels are exploratory and this remains
the principal internal-validity gap.
\fi

\section{Known Scope Boundaries and Remaining Gaps}
\label{sec:gaps}

Collected in one place so a reviewer need not reconstruct them.

\begin{enumerate}
\item \textbf{The effect/capability ontology is not populated by any extractor,
and so has no per-framework extraction mapping}
(Section~\ref{sec:ontology-gap}). The fields exist in the model and are populated
only when the curated corpus is authored, from declared tool names, by a
framework-independent keyword rule. \slug{state\_update} is never assigned by any
code path. \slug{back\_edge} is set only by \artifact{build\_corpus.py}.
\emph{The main paper's Section~2 now states this explicitly, and the two
documents agree}: what Table~\ref{tab:fwmap} documents is the per-framework
mapping into the exclusive \slug{NodeKind}/\slug{EdgeKind} view, which is the
only one that runs during extraction.
\item \textbf{Event-trace conservatism (H1) is assumed, not verified}
(Section~\ref{sec:h1}). Theorem~\ref{thm:etc} is conditional; provenance
metadata qualifies witnesses and only the caller's assertion discharges the hypothesis. A
body-level effect analysis is the missing piece.
\item \textbf{Per-framework extractor accuracy ``on stubs'' is not available}
(Section~\ref{sec:stubs}). \artifact{corpus/ground\_truth/} does not exist and the
available substitute is circular. Real-code accuracy is reported instead.
\item \textbf{No quantitative Promela/SMV modeling-effort measurement exists}
(Section~\ref{sec:eval}). One hand-written Promela model is committed; no timing,
no lines-of-model count, no head-to-head run.
\item \textbf{Confidence-output consistency has been repaired.}
The submitted \artifact{confidence\_intervals.json} is regenerated on the
matched, collision-free $n=106$ fidelity set with repository-clustered
bootstrap intervals and the self-repository-excluded 186-flag triage set. It
also records HGP-1 as $12/119$. The end-to-end consistency check rejects the
superseded $n=120/187$ output.
\item \textbf{HGP-1 supersedes the earlier human-gate estimand}
(Section~\ref{sec:hgp}). The primary figure is $12/119$; the earlier $3/119$ is
withdrawn, not averaged in. \emph{Residual gap:} HGP-1 is single-rater, 11 cases
remain \slug{arguable} for want of a repository-level corpus. The $12/119$
count is conditional on the upstream screen: the 87 workflows with no
source-level side-effecting tool were not re-read \emph{under HGP-1}---they were
excluded by the same
classification pass that produced the 32, and an independent effect-signature
sweep of all 87 flagged 10 for manual re-reading, of which 9 are correctly
screened and 1 would be \slug{arguable}. Screening error and label error can
move the observed count in either direction.
\item \textbf{Absolute benchmark timings are not comparable across drafts}
(Section~\ref{sec:eval}). Three recorded figures at 5{,}000 nodes span a factor
of ${\approx}33$ across machines. Only the asymptotic shape is a property of the
system.
\item \textbf{Independent human-validation status: \humanValidationStatus}
(Section~\ref{sec:human-validation}). Reference graphs and triage labels
originate in LLM passes sharing a model family. The survival test of
Section~\ref{sec:survival} weakens, but does not by itself remove, this threat:
it reproduces the headline decomposition from graphs alone without consulting
the labels ($93.0\%$ agreement), which is evidence against label
\emph{dependence}, not evidence of label \emph{correctness}.
\item \textbf{The reference graphs are not a semantic oracle, and we now
measure by how much} (Sections~\ref{sec:hgp},~\ref{sec:survival}). Node kinds
are assigned from naming and framework type, so a \slug{human} node need not be
able to block; $5/186$ flags trip the validity guard for this reason, and $9$ of
the $12$ source-audited violations are invisible even to the reconstruction because
the effect lives in a node body. Every human-gate estimand computed over
reference graphs is therefore biased downward by an amount we bound but cannot
correct.
\end{enumerate}

\bibliographystyle{plainnat}
\bibliography{../references}

\end{document}
